% ============================================================================
% MÓDULO 5 - Bases de datos, SQL y ETL desde R
% ============================================================================
\chapter{Bases de datos, SQL y ETL desde R}
\label{cap:bases-datos}

\begin{objetivos}
\begin{itemize}
  \item Explicar por qué las empresas guardan sus datos en bases de datos y
        no en archivos de Excel o CSV, y distinguir una base transaccional
        (OLTP) de una analítica (OLAP o \emph{data warehouse}).
  \item Usar el vocabulario básico: tabla, fila, columna, tipo de dato,
        llave primaria, llave foránea e índice.
  \item Conectarte desde R a una base de datos con \paquete{DBI} y
        \paquete{RSQLite}, cargar tablas, consultarlas y cerrar la conexión
        de forma segura.
  \item Escribir consultas SQL desde cero para responder preguntas de
        negocio: filtros, agregaciones, \emph{joins}, subconsultas, CTE y
        funciones de ventana.
  \item Proteger tus consultas contra la inyección SQL con consultas
        parametrizadas.
  \item Trabajar con tablas de la base como si fueran data frames gracias a
        \paquete{dbplyr}, y ver el SQL que genera.
  \item Diseñar y construir un \emph{data warehouse} en estrella con tabla
        de hechos, dimensiones, vistas e índices.
  \item Programar un proceso ETL reutilizable con transacciones, carga
        incremental sin duplicados y bitácora de ejecución.
  \item Conectarte de forma segura a las bases de datos reales de tu
        empresa (SQL Server, PostgreSQL, MySQL) sin exponer contraseñas.
\end{itemize}
\end{objetivos}

Hasta ahora has trabajado con archivos CSV que viven en tu computadora. En
una empresa real las cosas son distintas. Las ventas no están en un CSV: se
registran, una por una, en el sistema del punto de venta de cada sucursal, y
ese sistema las guarda en una \termino{base de datos}. El catálogo de
productos lo mantiene el área de compras en otro sistema; los clientes, el
CRM; los empleados, Recursos Humanos. Cuando tu jefa te pide ``las ventas de
ayer por sucursal'', nadie te va a mandar un archivo: te van a dar un usuario
y una contraseña para que consultes tú mismo la base de datos.

Por eso este módulo es clave en tu camino como analista de BI. Aprenderás
\termino{SQL} (\emph{Structured Query Language}), el idioma universal para
pedirle datos a una base; lo usan prácticamente todas las empresas del mundo
y aparece en casi cualquier vacante de analista de datos. Lo mejor: lo vas a
aprender \emph{desde R}, sin salir de RStudio, de modo que puedes combinar la
potencia de SQL para extraer y agregar datos con todo lo que ya sabes de
\paquete{dplyr} y \paquete{ggplot2}.

Para que puedas practicar sin instalar ningún servidor usaremos
\termino{SQLite}, un motor de base de datos completo que guarda toda la base
en un solo archivo. Construiremos, paso a paso, la base
\archivo{datos/bi\_retail.sqlite} con los datos del curso; primero como una
copia de los CSV, después como un \emph{data warehouse} en estrella cargado
por un proceso ETL escrito en R. Todo lo que aprendas aquí (el SQL, las
funciones de \paquete{DBI}, el diseño del modelo) funciona igual con SQL
Server, PostgreSQL o MySQL; al final verás cómo conectarte a ellos.

% ----------------------------------------------------------------------------
\section{¿Por qué bases de datos en BI?}
\label{sec:m5-porque}

\subsection{Los límites de Excel y los archivos CSV}

Excel y los CSV son excelentes para intercambiar datos pequeños, pero se
quedan cortos en cuanto una empresa crece:

\begin{itemize}
  \item \textbf{Volumen.} Una hoja de Excel admite 1\,048\,576 filas. Una
        cadena de 10 tiendas con 2\,000 tickets diarios por tienda llena ese
        límite en menos de dos meses. Un CSV no tiene límite formal, pero
        para leerlo tienes que cargarlo completo en la memoria de tu
        computadora.
  \item \textbf{Una sola versión de la verdad.} Cuando cada área guarda su
        propio archivo (\archivo{ventas\_final.xlsx},
        \archivo{ventas\_final\_v2\_bueno.xlsx}\ldots), en la junta cada quien
        llega con una cifra distinta. Una base de datos central es la
        \termino{fuente única de la verdad}: todos consultan los mismos
        datos.
  \item \textbf{Concurrencia.} Cien cajeros pueden registrar ventas al mismo
        tiempo en una base de datos sin pisarse; en un archivo de Excel
        compartido, el segundo que lo abre lo ve ``de solo lectura''.
  \item \textbf{Integridad.} Una base de datos puede impedir que se registre
        una venta de un producto que no existe o con cantidad negativa. Un
        CSV acepta cualquier cosa.
  \item \textbf{Seguridad y auditoría.} En una base puedes dar permiso de
        solo lectura a los analistas, ocultar columnas sensibles (sueldos,
        teléfonos) y saber quién consultó qué.
\end{itemize}

\subsection{Bases transaccionales (OLTP) y analíticas (OLAP)}

No todas las bases de datos tienen el mismo propósito. En BI conviene
distinguir dos grandes familias:

\begin{concepto}{OLTP frente a OLAP}
Una base \termino{OLTP} (\emph{Online Transaction Processing}) es la que
usan los sistemas del día a día: el punto de venta, el ERP, la tienda en
línea. Está optimizada para registrar muchísimas operaciones pequeñas y
rápidas (``inserta este ticket'', ``actualiza el stock de este producto'').

Una base \termino{OLAP} (\emph{Online Analytical Processing}), cuya forma
más común es el \termino{data warehouse} (almacén de datos), está pensada
para \emph{analizar}: pocas consultas, pero cada una recorre millones de
filas para sumar, agrupar y comparar (``ventas por mes, categoría y zona de
los últimos tres años''). Se alimenta de las bases OLTP mediante procesos
\termino{ETL} (\emph{Extract, Transform, Load}: extraer, transformar,
cargar).
\end{concepto}

\begin{table}[H]
\centering
\small
\caption{Diferencias entre una base transaccional y un data warehouse.}
\label{tab:m5-oltp-olap}
\begin{tabularx}{\textwidth}{@{}lXX@{}}
\toprule
 & \textbf{OLTP (transaccional)} & \textbf{OLAP (data warehouse)} \\
\midrule
Propósito & Operar el negocio & Analizar el negocio \\
Usuarios & Cajeros, sistemas, clientes & Analistas, directivos, tableros \\
Operación típica & Insertar o actualizar una fila & Leer y agregar millones de filas \\
Diseño & Muy normalizado (muchas tablas) & En estrella (hechos y dimensiones) \\
Historia & Estado actual & Años de historia \\
Ejemplo & Sistema del punto de venta & Base que alimenta el tablero de ventas \\
\bottomrule
\end{tabularx}
\end{table}

\begin{negocio}{¿Dónde encaja el analista de BI?}
Casi nunca vas a escribir directamente en la base OLTP de la empresa (y no
deberías: una consulta pesada puede volver lento el punto de venta). Tu
trabajo habitual es \emph{leer} de las bases con permisos de solo lectura,
construir o alimentar el \emph{data warehouse} y, a partir de él, producir
reportes, tableros y modelos. R es una herramienta excelente para las tres
tareas.
\end{negocio}

\subsection{Motores de bases de datos que encontrarás en las empresas}

Un \termino{motor de base de datos} (o \termino{SGBD}, sistema gestor de
bases de datos) es el programa que guarda los datos y responde las
consultas. Todos hablan SQL, con pequeñas diferencias de ``dialecto''. La
tabla~\ref{tab:m5-motores} resume los más comunes y el paquete de R con el
que te conectas a cada uno.

\begin{table}[H]
\centering
\small
\caption{Motores de bases de datos frecuentes y cómo conectarte desde R.}
\label{tab:m5-motores}
\begin{tabularx}{\textwidth}{@{}lXl@{}}
\toprule
\textbf{Motor} & \textbf{Dónde lo verás} & \textbf{Paquete de R} \\
\midrule
SQL Server & Empresas con ecosistema Microsoft, ERPs & \paquete{odbc} \\
PostgreSQL & Startups, aplicaciones web, gobierno & \paquete{RPostgres} \\
MySQL / MariaDB & Tiendas en línea, sistemas web & \paquete{RMariaDB} \\
Oracle & Bancos, telecomunicaciones, grandes corporativos & \paquete{odbc} \\
SQLite & Aplicaciones móviles, prototipos, este curso & \paquete{RSQLite} \\
BigQuery, Snowflake, Redshift & Data warehouses en la nube & \paquete{bigrquery}, \paquete{odbc} \\
\bottomrule
\end{tabularx}
\end{table}

\begin{consejo}
Lo que aprendas con SQLite te sirve para cualquier motor: las funciones de
\paquete{DBI} son idénticas y el 90\% del SQL también. Las diferencias
suelen estar en funciones de fecha (\codigo{strftime()} en SQLite,
\codigo{DATEPART()} en SQL Server, \codigo{date\_trunc()} en PostgreSQL) y en
cómo limitar filas (\codigo{LIMIT} frente a \codigo{TOP}).
\end{consejo}

% ----------------------------------------------------------------------------
\section{Conceptos básicos de una base de datos relacional}
\label{sec:m5-conceptos}

Una \termino{base de datos relacional} organiza la información en
\termino{tablas} relacionadas entre sí. Si ya dominas los data frames, vas
con ventaja: una tabla se parece mucho a un data frame.

\begin{concepto}{Vocabulario esencial}
\begin{itemize}
  \item \textbf{Tabla}: conjunto de datos sobre un mismo tipo de cosa
        (clientes, productos, ventas). Equivale a un data frame.
  \item \textbf{Fila} (registro): un elemento concreto, por ejemplo un
        cliente.
  \item \textbf{Columna} (campo): un atributo, por ejemplo la ciudad del
        cliente. Cada columna tiene un \termino{tipo de dato}: entero
        (\codigo{INTEGER}), decimal (\codigo{REAL}), texto (\codigo{TEXT}),
        fecha, etc.
  \item \textbf{Llave primaria} (\emph{primary key}, PK): columna cuyo valor
        identifica de forma única cada fila; no se repite ni puede estar
        vacía. Ejemplo: \codigo{cliente\_id}.
  \item \textbf{Llave foránea} (\emph{foreign key}, FK): columna de una
        tabla que apunta a la llave primaria de otra. En
        \codigo{transacciones}, la columna \codigo{cliente\_id} es una llave
        foránea hacia \codigo{clientes}. Así se \emph{relacionan} las tablas.
  \item \textbf{Índice}: estructura auxiliar (como el índice alfabético de
        un libro) que permite encontrar filas rápidamente sin recorrer toda
        la tabla.
\end{itemize}
\end{concepto}

La figura~\ref{fig:m5-relacional} muestra el modelo relacional de los datos
del curso. La tabla central es \codigo{transacciones}: cada fila es una venta
y, mediante sus llaves foráneas, sabemos a qué cliente, producto, tienda y
vendedor corresponde. Observa que \codigo{vendedor\_id} apunta a
\codigo{empleado\_id} en la tabla de empleados: la columna no tiene que
llamarse igual en ambas tablas.

\begin{figure}[H]
\centering
\begin{tikzpicture}[
  tabla/.style={draw=primario, rounded corners=2pt, fill=white,
    font=\ttfamily\scriptsize, align=left, inner sep=3pt,
    minimum width=3.3cm},
  titulo/.style={fill=primario, text=white, font=\sffamily\small\bfseries,
    minimum width=3.3cm, inner sep=3pt},
  rel/.style={-{Stealth[length=2.5mm]}, draw=acento, thick}
]
% Tabla central
\node[titulo] (tt) at (0,0) {transacciones};
\node[tabla, below=0pt of tt] (t) {\textbf{PK} transaccion\_id\\
fecha, hora\\ \textbf{FK} cliente\_id\\ \textbf{FK} producto\_id\\
\textbf{FK} tienda\_id\\ \textbf{FK} vendedor\_id\\
cantidad, precio\_venta\\ metodo\_pago\\ total\_transaccion};
% Clientes
\node[titulo] (ct) at (-5.2,1.2) {clientes};
\node[tabla, below=0pt of ct] (c) {\textbf{PK} cliente\_id\\ nombre, edad\\
ciudad, segmento\\ es\_premium};
% Productos
\node[titulo] (pt) at (5.2,1.2) {productos};
\node[tabla, below=0pt of pt] (p) {\textbf{PK} producto\_id\\
nombre\_producto\\ categoria, marca\\ precio\_catalogo, costo};
% Tiendas
\node[titulo] (tit) at (-5.2,-2.6) {tiendas};
\node[tabla, below=0pt of tit] (ti) {\textbf{PK} tienda\_id\\
nombre\_tienda\\ ciudad, zona, tipo};
% Empleados
\node[titulo] (et) at (5.2,-2.6) {empleados};
\node[tabla, below=0pt of et] (e) {\textbf{PK} empleado\_id\\ nombre\\
departamento, puesto};
% Relaciones (de la FK a la PK)
\draw[rel] ($(t.west)+(0,0.9)$) -- node[above, sloped,
  font=\scriptsize\sffamily, text=gristexto] {cliente\_id} (c.east);
\draw[rel] ($(t.east)+(0,0.9)$) -- node[above, sloped,
  font=\scriptsize\sffamily, text=gristexto] {producto\_id} (p.west);
\draw[rel] ($(t.west)+(0,-0.6)$) -- node[below, sloped,
  font=\scriptsize\sffamily, text=gristexto] {tienda\_id} (ti.east);
\draw[rel] ($(t.east)+(0,-0.6)$) -- node[below, sloped,
  font=\scriptsize\sffamily, text=gristexto] {vendedor\_id} (e.west);
\end{tikzpicture}
\caption{Modelo relacional de los datos del curso. Las flechas van de la
llave foránea (FK) a la llave primaria (PK) de la tabla relacionada.}
\label{fig:m5-relacional}
\end{figure}

\begin{advertencia}
En los datos del curso las llaves foráneas existen ``de palabra'': los CSV
no obligan a nada. Una base de datos bien diseñada \emph{declara} esas
relaciones para que el motor rechace, por ejemplo, una venta de un producto
inexistente. Lo harás en la sección~\ref{sec:m5-estrella}, cuando construyas
el \emph{data warehouse}.
\end{advertencia}

\begin{hazlotu}
Abre \archivo{ventas\_retail.csv} y responde en papel: ¿cuál sería su llave
primaria? (Pista: no tiene una columna \codigo{\_id} propia, pero hay una
venta por día.) ¿Qué columnas son llaves foráneas y hacia qué tablas
apuntan?
\end{hazlotu}

% ----------------------------------------------------------------------------
\section{Conectarte a una base de datos con DBI y RSQLite}
\label{sec:m5-dbi}

En R, todas las conexiones a bases de datos siguen el mismo esquema de dos
capas:

\begin{itemize}
  \item \paquete{DBI} (\emph{DataBase Interface}) define las funciones que
        tú usas: \codigo{dbConnect()}, \codigo{dbGetQuery()},
        \codigo{dbWriteTable()}, etc. Son las mismas para cualquier motor.
  \item Un \termino{driver} (controlador) traduce esas funciones al
        protocolo de un motor concreto: \paquete{RSQLite} para SQLite,
        \paquete{RPostgres} para PostgreSQL, \paquete{odbc} para SQL Server,
        etc.
\end{itemize}

Si mañana tu empresa cambia de SQLite a PostgreSQL, solo cambia la línea de
\codigo{dbConnect()}; el resto de tu código sigue igual. Los paquetes que
usaremos se instalan como siempre:

\begin{rnoejecutar}
install.packages(c("DBI", "RSQLite", "dbplyr"))
\end{rnoejecutar}

\subsection{Preparar la carpeta y crear la base}

Empezamos cargando los paquetes y preparando la ruta del archivo de la base.
Un detalle importante de reproducibilidad: SQLite \emph{guarda} lo que
escribes en el archivo, y ese archivo sigue ahí la próxima vez que ejecutes
el script. Si lo corres dos veces, la segunda vez las tablas ya existirían y
obtendrías errores del tipo ``la tabla ya existe'' o, peor, datos
duplicados. Por eso, al inicio del módulo borramos el archivo si existe y
empezamos siempre desde cero.

\begin{rcode}
library(tidyverse)
library(DBI)        # funciones comunes para cualquier base de datos
library(RSQLite)    # driver (controlador) de SQLite
library(dbplyr)     # traduce verbos de dplyr a SQL
library(lubridate)  # fechas
\end{rcode}

Al cargar \paquete{dbplyr} verás un aviso de que ``enmascara'' las
funciones \codigo{ident()} y \codigo{sql()} de \paquete{dplyr}; es normal y
no afecta nada de lo que haremos. Ahora preparamos la carpeta y la ruta del
archivo:

\begin{rcode}
# Carpeta donde vivirá la base de datos
dir.create("datos", showWarnings = FALSE)
ruta_bd <- "datos/bi_retail.sqlite"

# Si la base ya existe (de una ejecución anterior), la borramos para
# que el script sea reproducible y siempre empiece desde cero.
# unlink() no marca error si el archivo no existe.
unlink(ruta_bd)
file.exists(ruta_bd)
\end{rcode}

\begin{rsalida}
[1] FALSE
\end{rsalida}

Ahora abrimos la conexión. Con SQLite, si el archivo no existe,
\codigo{dbConnect()} lo crea vacío:

\begin{rcode}
# Abrir (y crear) la base de datos en el archivo
con <- dbConnect(RSQLite::SQLite(), ruta_bd)

# ¿Qué tablas tiene? Ninguna todavía
dbListTables(con)
file.exists(ruta_bd)
\end{rcode}

\begin{rsalida}
character(0)
[1] TRUE
\end{rsalida}

El objeto \codigo{con} es la \termino{conexión}: el ``cable'' entre tu
sesión de R y la base. Todas las funciones de \paquete{DBI} reciben la
conexión como primer argumento.

\begin{consejo}
SQLite también puede crear una base \emph{temporal en memoria} con
\codigo{dbConnect(RSQLite::SQLite(), ":memory:")}. Desaparece al
desconectarte, así que es perfecta para pruebas rápidas y ejercicios; las
soluciones de este módulo la usan mucho.
\end{consejo}

\subsection{Cargar los CSV del curso a la base con \texorpdfstring{\codigo{dbWriteTable()}}{dbWriteTable()}}

\codigo{dbWriteTable(con, "nombre", datos)} crea una tabla en la base a
partir de un data frame. Vamos a cargar los seis CSV del curso. Antes de
escribirlos haremos una pequeña transformación: convertir las columnas de
fecha y de hora a texto con formato ISO (\codigo{"2023-01-13"},
\codigo{"13:15:00"}). SQLite no tiene un tipo de dato de fecha propiamente
dicho y, si le mandas un \codigo{Date} de R tal cual, lo guarda como un
número. En la sección~\ref{sec:m5-fechas} verás el problema con detalle; por
ahora quédate con la regla: \emph{en SQLite, las fechas se guardan como texto
\codigo{'AAAA-MM-DD'}}.

\begin{rcode}
# Lee un CSV del curso y deja fechas y horas como texto ISO
leer_para_bd <- function(archivo) {
  read_csv(file.path("datasets", archivo), show_col_types = FALSE) %>%
    mutate(
      across(where(is.Date), ~ format(.x, "%Y-%m-%d")),  # fechas
      across(where(hms::is_hms), as.character)          # horas
    )
}

# Nombre de la tabla en la base = nombre del archivo sin ".csv"
archivos <- c("ventas_retail", "clientes", "productos",
              "empleados", "transacciones", "tiendas")

for (tabla in archivos) {
  datos <- leer_para_bd(paste0(tabla, ".csv"))
  # overwrite = TRUE: si la tabla existe, la reemplaza
  dbWriteTable(con, tabla, datos, overwrite = TRUE)
  cat(sprintf("%-14s %5d filas cargadas\n", tabla, nrow(datos)))
}

dbListTables(con)
\end{rcode}

\begin{rsalida}
ventas_retail    365 filas cargadas
clientes         200 filas cargadas
productos         50 filas cargadas
empleados        100 filas cargadas
transacciones   1000 filas cargadas
tiendas           10 filas cargadas
[1] "clientes"      "empleados"     "productos"     "tiendas"      
[5] "transacciones" "ventas_retail"
\end{rsalida}

En unas cuantas líneas tienes una base de datos con seis tablas.
\codigo{dbListTables()} las devuelve en orden alfabético. Si abres la
carpeta \archivo{datos/} verás el archivo \archivo{bi\_retail.sqlite}: esa es
toda la base; puedes copiarlo, respaldarlo o mandarlo por correo.

\subsection{Explorar la base: columnas, tipos y contenido}

\codigo{dbListFields()} te dice qué columnas tiene una tabla, y la
instrucción especial de SQLite \codigo{PRAGMA table\_info()} te muestra
además el tipo de dato con el que se guardó cada una:

\begin{rcode}
# Columnas de la tabla de tiendas
dbListFields(con, "tiendas")

# Tipo de dato de cada columna (instrucción propia de SQLite)
dbGetQuery(con, "PRAGMA table_info(tiendas)") %>%
  select(name, type)
\end{rcode}

\begin{rsalida}
[1] "tienda_id"      "nombre_tienda"  "ciudad"        
[4] "zona"           "tamano_m2"      "empleados"     
[7] "fecha_apertura" "tipo"          
            name type
1      tienda_id REAL
2  nombre_tienda TEXT
3         ciudad TEXT
4           zona TEXT
5      tamano_m2 REAL
6      empleados REAL
7 fecha_apertura TEXT
8           tipo TEXT
\end{rsalida}

Observa la traducción de tipos: los números de R (\codigo{double}) se
guardaron como \codigo{REAL}, el texto como \codigo{TEXT} y la fecha de
apertura, gracias a nuestra conversión, también como \codigo{TEXT}. Los
valores lógicos (\codigo{TRUE}/\codigo{FALSE}, como \codigo{es\_premium} en
clientes) se guardan como enteros \codigo{1}/\codigo{0}, porque SQLite no
tiene tipo lógico.

Para traer una tabla \emph{completa} a R usa \codigo{dbReadTable()}. Devuelve
un data frame normal:

\begin{rcode}
tiendas_bd <- dbReadTable(con, "tiendas")
class(tiendas_bd)
tiendas_bd %>% select(tienda_id, nombre_tienda, ciudad, zona) %>% head(4)
\end{rcode}

\begin{rsalida}
[1] "data.frame"
  tienda_id   nombre_tienda      ciudad   zona
1         1 Sucursal Centro        CDMX Centro
2         2  Sucursal Norte Guadalajara  Norte
3         3    Sucursal Sur   Monterrey  Norte
4         4   Sucursal Este        CDMX   Este
\end{rsalida}

\begin{advertencia}
\codigo{dbReadTable()} trae \emph{todas} las filas y columnas. Con la tabla
de tiendas (10 filas) no pasa nada, pero con la tabla de ventas de una
empresa real (cientos de millones de filas) puedes congelar tu computadora.
Lo profesional es pedirle a la base \emph{solo lo que necesitas} con una
consulta SQL, que es justo lo que viene.
\end{advertencia}

\subsection{Consultar y modificar: \texorpdfstring{\codigo{dbGetQuery()}}{dbGetQuery()} y \texorpdfstring{\codigo{dbExecute()}}{dbExecute()}}

Las dos funciones que más vas a usar son:

\begin{itemize}
  \item \codigo{dbGetQuery(con, sql)}: ejecuta una consulta que
        \emph{devuelve datos} (\codigo{SELECT}) y te regresa un data frame.
  \item \codigo{dbExecute(con, sql)}: ejecuta una instrucción que
        \emph{modifica} la base (\codigo{CREATE}, \codigo{INSERT},
        \codigo{UPDATE}, \codigo{DELETE}) y te regresa el número de filas
        afectadas.
\end{itemize}

\begin{rcode}
# Una consulta: ¿cuántas transacciones hay?
dbGetQuery(con, "SELECT COUNT(*) AS num_transacciones
                 FROM transacciones")

# Una modificación: crear una tabla de metas mensuales y llenarla
dbExecute(con, "CREATE TABLE metas (mes TEXT, meta REAL)")
dbExecute(con, "INSERT INTO metas VALUES ('2023-01', 80000),
                                         ('2023-02', 75000)")
dbGetQuery(con, "SELECT * FROM metas")

# Borrar la tabla de prueba
dbExecute(con, "DROP TABLE metas")
\end{rcode}

\begin{rsalida}
  num_transacciones
1              1000
[1] 0
[1] 2
      mes  meta
1 2023-01 80000
2 2023-02 75000
[1] 0
\end{rsalida}

El texto SQL puede ocupar varias líneas dentro de las comillas; a la base no
le importan los saltos de línea. \codigo{CREATE TABLE} devolvió \codigo{0}
(no afectó filas), el \codigo{INSERT} devolvió \codigo{2} (insertó dos filas)
y el \codigo{DROP TABLE} volvió a devolver \codigo{0}.

\subsection{Cerrar la conexión: \texorpdfstring{\codigo{dbDisconnect()}}{dbDisconnect()} y \texorpdfstring{\codigo{on.exit()}}{on.exit()}}

Cuando terminas, cierra la conexión con \codigo{dbDisconnect(con)}. Con
SQLite no cerrarla es un descuido menor, pero en un servidor real cada
conexión abierta consume recursos y muchas empresas limitan cuántas puede
tener abierta cada usuario. Si escribes una \emph{función} que abre su
propia conexión, la buena práctica es usar \codigo{on.exit()}: registra una
acción que se ejecutará al salir de la función, \emph{incluso si ocurre un
error a la mitad}.

\begin{rcode}
# Función que abre su propia conexión y SIEMPRE la cierra
contar_filas <- function(ruta, tabla) {
  con_local <- dbConnect(RSQLite::SQLite(), ruta)
  on.exit(dbDisconnect(con_local))   # se ejecuta al salir, pase lo que pase
  sql <- paste("SELECT COUNT(*) AS n FROM", tabla)
  dbGetQuery(con_local, sql)$n
}

contar_filas(ruta_bd, "clientes")
contar_filas(ruta_bd, "productos")
\end{rcode}

\begin{rsalida}
[1] 200
[1] 50
\end{rsalida}

Aquí usamos \codigo{paste()} para armar el SQL porque el nombre de la tabla
lo escribes tú, en tu propio código. En la sección~\ref{sec:m5-parametros}
verás por qué \emph{nunca} debes hacer lo mismo con valores que escribe un
usuario. En el resto del módulo seguiremos usando la conexión global
\codigo{con}, que cerraremos al final.

\begin{hazlotu}
Usa \codigo{dbListFields()} para ver las columnas de la tabla
\codigo{empleados} y \codigo{contar\_filas()} para confirmar que tiene 100
filas.
\end{hazlotu}

% ----------------------------------------------------------------------------
\section{SQL desde cero con preguntas de negocio}
\label{sec:m5-sql}

SQL es un lenguaje \emph{declarativo}: le dices a la base \emph{qué} quieres,
no \emph{cómo} conseguirlo; el motor decide la forma más eficiente de
obtenerlo. Una consulta tiene siempre la misma estructura, con las
cláusulas en este orden (las que están entre corchetes son opcionales):

\begin{sqlcode}
SELECT   columnas o cálculos          -- ¿qué quiero ver?
FROM     tabla                        -- ¿de dónde?
[JOIN    otra_tabla ON condición]     -- ¿con qué la combino?
[WHERE   condición sobre filas]       -- ¿qué filas me interesan?
[GROUP BY columnas]                   -- ¿por qué agrupo?
[HAVING  condición sobre grupos]      -- ¿qué grupos me interesan?
[ORDER BY columnas]                   -- ¿en qué orden?
[LIMIT   n]                           -- ¿cuántas filas como máximo?
\end{sqlcode}

Algunas convenciones: las palabras clave de SQL se escriben en mayúsculas
por costumbre (no es obligatorio), los comentarios empiezan con
\codigo{-{}-}, el texto va entre \emph{comillas simples}
(\codigo{'CDMX'}) y cada consulta puede terminar con punto y coma. En R,
toda la consulta va dentro de una cadena de texto que le pasas a
\codigo{dbGetQuery()}.

\begin{concepto}{SQL y dplyr piensan igual}
Si dominas \paquete{dplyr} (Módulo~\ref{cap:manipulacion}), ya conoces la
lógica de SQL: \codigo{SELECT} es \codigo{select()}/\codigo{mutate()},
\codigo{WHERE} es \codigo{filter()}, \codigo{GROUP BY} +
funciones de agregación es \codigo{group\_by()} + \codigo{summarise()},
\codigo{ORDER BY} es \codigo{arrange()} y \codigo{JOIN} es
\codigo{left\_join()} y compañía. De hecho, \paquete{dplyr} se diseñó
inspirándose en SQL. Al final del módulo tienes una tabla de equivalencias.
\end{concepto}

\subsection{SELECT y FROM: elegir columnas}

\begin{ejemplo}{Directorio de sucursales}
Recursos Humanos quiere un directorio con el nombre, la ciudad y la zona de
cada sucursal.
\end{ejemplo}

\begin{rcode}
dbGetQuery(con, "
  SELECT nombre_tienda, ciudad, zona   -- columnas que quiero
  FROM tiendas                         -- tabla de origen
")
\end{rcode}

\begin{rsalida}
      nombre_tienda      ciudad   zona
1   Sucursal Centro        CDMX Centro
2    Sucursal Norte Guadalajara  Norte
3      Sucursal Sur   Monterrey  Norte
4     Sucursal Este        CDMX   Este
5    Sucursal Oeste      Puebla Centro
6  Sucursal Plaza A        CDMX  Oeste
7  Sucursal Plaza B Guadalajara  Norte
8   Sucursal Outlet     Tijuana  Norte
9  Sucursal Premium        CDMX    Sur
10 Sucursal Express   Querétaro Centro
\end{rsalida}

\codigo{SELECT *} trae todas las columnas; es cómodo para explorar, pero en
reportes conviene listar solo las que necesitas. \codigo{LIMIT} restringe el
número de filas, ideal para ``echar un vistazo'' a una tabla grande:

\begin{rcode}
dbGetQuery(con, "SELECT * FROM transacciones LIMIT 3")
\end{rcode}

\begin{rsalida}
  transaccion_id      fecha     hora cliente_id producto_id cantidad
1              1 2023-01-13 13:15:00        155           3        4
2              2 2023-01-09 10:45:00        155          47        4
3              3 2023-02-20 10:30:00         75           6        3
  precio_venta     metodo_pago tienda_id vendedor_id
1       303.97 Tarjeta Crédito        10           4
2        58.07  Tarjeta Débito         2          22
3       275.46  Tarjeta Débito         7          14
  total_transaccion     mes trimestre dia_semana
1           1215.88 2023-01        Q1    viernes
2            232.28 2023-01        Q1      lunes
3            826.38 2023-02        Q1      lunes
\end{rsalida}

Como la tabla tiene muchas columnas, R la imprime en varios renglones. Nota
que la fecha y la hora se ven como texto, justo como las guardamos.

\subsection{WHERE: filtrar filas}

\codigo{WHERE} se queda solo con las filas que cumplen una condición. Los
operadores de comparación son \codigo{=} (¡un solo signo igual, no dos como
en R!), \codigo{<>} o \codigo{!=} (distinto), \codigo{<}, \codigo{>},
\codigo{<=} y \codigo{>=}. Las condiciones se combinan con \codigo{AND},
\codigo{OR} y \codigo{NOT}.

\begin{ejemplo}{Clientes premium del norte}
Marketing prepara una campaña para los clientes premium de Monterrey y
Guadalajara. ¿Cuántos son y quiénes?
\end{ejemplo}

\begin{rcode}
# IN (...) equivale a varias condiciones unidas con OR
dbGetQuery(con, "
  SELECT COUNT(*) AS clientes
  FROM clientes
  WHERE es_premium = 1
    AND ciudad IN ('Monterrey', 'Guadalajara')
")

dbGetQuery(con, "
  SELECT nombre, ciudad, edad
  FROM clientes
  WHERE es_premium = 1
    AND ciudad IN ('Monterrey', 'Guadalajara')
  ORDER BY ciudad, nombre
  LIMIT 6
")
\end{rcode}

\begin{rsalida}
  clientes
1       16
              nombre      ciudad edad
1       Carmen Reyes Guadalajara   33
2 Fernando Hernández Guadalajara   26
3     Marta González Guadalajara   38
4      Marta Morales Guadalajara   55
5        Miguel Díaz Guadalajara   62
6     Ricardo Torres Guadalajara   72
\end{rsalida}

Recuerda que los lógicos se guardaron como \codigo{1}/\codigo{0}, por eso
filtramos con \codigo{es\_premium = 1}. La primera consulta cuenta, la
segunda muestra los primeros seis ordenados por ciudad y nombre.

\begin{ejemplo}{Ventas navideñas y tickets grandes}
Queremos ver las transacciones entre el 22 y el 24 de diciembre y saber
cuántos tickets superan los \$2\,800.
\end{ejemplo}

\begin{rcode}
# BETWEEN incluye ambos extremos; las fechas son texto ISO
dbGetQuery(con, "
  SELECT transaccion_id, fecha, producto_id, total_transaccion
  FROM transacciones
  WHERE fecha BETWEEN '2023-12-22' AND '2023-12-24'
  ORDER BY fecha
")

dbGetQuery(con, "
  SELECT COUNT(*) AS tickets_grandes
  FROM transacciones
  WHERE total_transaccion > 2800
")
\end{rcode}

\begin{rsalida}
  transaccion_id      fecha producto_id total_transaccion
1            693 2023-12-22          17           1130.45
2            828 2023-12-22          32           1300.50
3            452 2023-12-23          31            480.65
4            558 2023-12-23          29            334.40
5            926 2023-12-23          27           2709.85
6            518 2023-12-24           3            356.36
7            796 2023-12-24          34            584.71
8            941 2023-12-24          17           1717.96
  tickets_grandes
1              26
\end{rsalida}

Comparar fechas con \codigo{BETWEEN} funciona porque el texto
\codigo{'AAAA-MM-DD'} se ordena alfabéticamente igual que cronológicamente
(\codigo{'2023-12-22'} va antes que \codigo{'2023-12-24'}). Esa es otra
razón para guardar las fechas en formato ISO.

\codigo{LIKE} busca patrones en texto: \codigo{\%} representa ``cualquier
cantidad de caracteres'' y \codigo{\_} ``exactamente un carácter''.

\begin{rcode}
# Clientes cuyo apellido es García y cuyo nombre empieza con "Laura"
dbGetQuery(con, "
  SELECT
    SUM(nombre LIKE '%García') AS apellido_garcia,
    SUM(nombre LIKE 'Laura%')  AS nombre_laura
  FROM clientes
")
\end{rcode}

\begin{rsalida}
  apellido_garcia nombre_laura
1              10            9
\end{rsalida}

Aquí usamos un truco útil: una comparación en SQLite vale \codigo{1} si es
verdadera y \codigo{0} si es falsa, así que \codigo{SUM()} de una condición
cuenta cuántas filas la cumplen (igual que \codigo{sum()} de un vector
lógico en R).

\begin{advertencia}
Cuando combines \codigo{AND} y \codigo{OR}, usa paréntesis.
\codigo{WHERE ciudad = 'CDMX' OR ciudad = 'Puebla' AND edad > 60} significa
``todos los de CDMX, más los de Puebla mayores de 60'', porque
\codigo{AND} se evalúa antes que \codigo{OR}. Si querías ``mayores de 60 de
CDMX o Puebla'', escribe
\codigo{WHERE (ciudad = 'CDMX' OR ciudad = 'Puebla') AND edad > 60}. Y para
valores vacíos no uses \codigo{= NULL}: usa \codigo{IS NULL} o
\codigo{IS NOT NULL}.
\end{advertencia}

\subsection{ORDER BY y LIMIT: los primeros de la lista}

\begin{ejemplo}{Los cinco tickets más grandes del año}
La gerencia quiere revisar las cinco ventas de mayor importe para validar
que no sean errores de captura.
\end{ejemplo}

\begin{rcode}
dbGetQuery(con, "
  SELECT transaccion_id, fecha, cantidad, precio_venta,
         total_transaccion
  FROM transacciones
  ORDER BY total_transaccion DESC   -- DESC = de mayor a menor
  LIMIT 5
")
\end{rcode}

\begin{rsalida}
  transaccion_id      fecha cantidad precio_venta total_transaccion
1            176 2023-04-03        5       598.47           2992.35
2            172 2023-06-14        5       598.09           2990.45
3            835 2023-08-19        5       596.74           2983.70
4            535 2023-06-25        5       595.24           2976.20
5            451 2023-05-02        5       594.71           2973.55
\end{rsalida}

\codigo{ORDER BY} ordena de menor a mayor por omisión (\codigo{ASC});
\codigo{DESC} invierte el orden. Los cinco tickets más grandes son de cinco
unidades con precio cercano al máximo de \$600: son plausibles, no parecen
errores de captura.

\subsection{Funciones de agregación, GROUP BY y HAVING}

Las \termino{funciones de agregación} resumen muchas filas en un solo
valor: \codigo{COUNT()}, \codigo{SUM()}, \codigo{AVG()}, \codigo{MIN()} y
\codigo{MAX()}. \codigo{COUNT(*)} cuenta filas y
\codigo{COUNT(DISTINCT columna)} cuenta valores distintos.

\begin{ejemplo}{KPIs generales del año}
Antes de cualquier análisis, el director quiere los números grandes de
2023: número de tickets, clientes distintos, ventas totales, ticket
promedio y ticket máximo.
\end{ejemplo}

\begin{rcode}
dbGetQuery(con, "
  SELECT COUNT(*)                         AS tickets,
         COUNT(DISTINCT cliente_id)       AS clientes,
         ROUND(SUM(total_transaccion), 2) AS ventas,
         ROUND(AVG(total_transaccion), 2) AS ticket_promedio,
         MAX(total_transaccion)           AS ticket_maximo
  FROM transacciones
")
\end{rcode}

\begin{rsalida}
  tickets clientes   ventas ticket_promedio ticket_maximo
1    1000      198 959992.7          959.99       2992.35
\end{rsalida}

Con \codigo{AS} pones un \termino{alias}, es decir, un nombre a la columna
calculada (sin él, la columna se llamaría literalmente
\codigo{ROUND(SUM(total\_transaccion), 2)}). En 2023 hubo 1\,000 tickets por
\$959\,992.7, con un ticket promedio de \$959.99. De los 200 clientes
registrados, 198 compraron al menos una vez.

\codigo{GROUP BY} calcula las agregaciones \emph{por grupo}, igual que
\codigo{group\_by()} + \codigo{summarise()}:

\begin{ejemplo}{¿Cómo pagan nuestros clientes?}
Finanzas negocia comisiones con el banco y necesita saber cuántas ventas y
cuánto dinero entra por cada método de pago.
\end{ejemplo}

\begin{rcode}
dbGetQuery(con, "
  SELECT metodo_pago,
         COUNT(*)                         AS tickets,
         ROUND(SUM(total_transaccion), 0) AS ventas,
         ROUND(AVG(total_transaccion), 2) AS ticket_promedio
  FROM transacciones
  GROUP BY metodo_pago
  ORDER BY ventas DESC
")
\end{rcode}

\begin{rsalida}
      metodo_pago tickets ventas ticket_promedio
1 Tarjeta Crédito     350 351171         1003.34
2        Efectivo     315 303735          964.24
3  Tarjeta Débito     233 220855          947.88
4   Transferencia     102  84232          825.80
\end{rsalida}

La tarjeta de crédito es el método más usado (350 tickets) y el que más
dinero trae (\$351\,171), además de tener el ticket promedio más alto
(\$1\,003). La transferencia es minoritaria: 102 tickets con un ticket
promedio de solo \$826. Si la comisión bancaria de la tarjeta de crédito es
alta, esta tabla dice cuánto dinero está en juego.

\codigo{WHERE} filtra \emph{filas antes} de agrupar; para filtrar
\emph{grupos después} de agregarlos se usa \codigo{HAVING}:

\begin{ejemplo}{Clientes frecuentes}
El programa de lealtad quiere invitar a los clientes con 10 compras o más en
el año.
\end{ejemplo}

\begin{rcode}
dbGetQuery(con, "
  SELECT cliente_id,
         COUNT(*)                         AS compras,
         ROUND(SUM(total_transaccion), 2) AS gasto_total
  FROM transacciones
  GROUP BY cliente_id
  HAVING COUNT(*) >= 10      -- condición sobre el grupo
  ORDER BY compras DESC, gasto_total DESC
")
\end{rcode}

\begin{rsalida}
  cliente_id compras gasto_total
1          5      13    14680.10
2         55      13    11355.70
3        105      13    10715.72
4        168      12    11675.62
5        125      10    11239.83
6        162      10    11176.45
\end{rsalida}

Solo seis clientes compraron 10 veces o más; el cliente 5 encabeza la lista
con 13 compras y \$14\,680. Observa que el \codigo{HAVING} usa
\codigo{COUNT(*)} (el valor del grupo), algo que \codigo{WHERE} no puede
hacer porque se evalúa antes de agrupar.

\begin{advertencia}
Toda columna del \codigo{SELECT} que no esté dentro de una función de
agregación debe aparecer en el \codigo{GROUP BY}. SQLite es permisivo y no
marca error si lo olvidas (devuelve un valor cualquiera del grupo), pero SQL
Server y PostgreSQL sí lo rechazan. Acostúmbrate a hacerlo bien desde el
principio.
\end{advertencia}

\subsection{CASE WHEN: clasificar dentro de la consulta}

\codigo{CASE WHEN} es el \codigo{case\_when()} de SQL: evalúa condiciones en
orden y devuelve el valor de la primera que se cumple.

\begin{ejemplo}{Distribución de tickets por rango}
El área comercial quiere saber cuántas ventas son pequeñas, medianas o
grandes, y cuánto aporta cada rango.
\end{ejemplo}

\begin{rcode}
dbGetQuery(con, "
  SELECT CASE
           WHEN total_transaccion < 500  THEN '1. Menos de 500'
           WHEN total_transaccion < 1500 THEN '2. De 500 a 1,500'
           ELSE '3. 1,500 o más'
         END AS rango_ticket,
         COUNT(*) AS tickets,
         ROUND(SUM(total_transaccion), 0) AS ventas
  FROM transacciones
  GROUP BY rango_ticket
  ORDER BY rango_ticket
")
\end{rcode}

\begin{rsalida}
       rango_ticket tickets ventas
1   1. Menos de 500     377 101783
2 2. De 500 a 1,500     379 347606
3    3. 1,500 o más     244 510603
\end{rsalida}

Los tickets de \$1\,500 o más son solo 244 de 1\,000 (menos de la cuarta
parte), pero aportan \$510\,603, más de la mitad de las ventas. Los 377
tickets pequeños apenas suman \$101\,783. Es la clásica regla de Pareto:
pocas ventas grandes explican gran parte del ingreso. Numerar las etiquetas (\codigo{'1. \ldots'}) es un
truco sencillo para que \codigo{ORDER BY} las deje en el orden lógico y no
en orden alfabético.

\begin{hazlotu}
Escribe una consulta que muestre, para cada ciudad, cuántos clientes hay y
su edad promedio, solo para las ciudades con más de 20 clientes, ordenadas
de mayor a menor número de clientes. (Pista: \codigo{GROUP BY ciudad} y
\codigo{HAVING COUNT(*) > 20}.)
\end{hazlotu}

\subsection{Fechas en SQLite}
\label{sec:m5-fechas}

Las fechas son la fuente número uno de dolores de cabeza al pasar datos de
R a SQLite. Veamos el problema con calma y luego la solución.

\subsubsection{El problema: los Date de R se guardan como números}

SQLite solo tiene cinco clases de almacenamiento: \codigo{NULL},
\codigo{INTEGER}, \codigo{REAL}, \codigo{TEXT} y \codigo{BLOB}; no existe un
tipo fecha. Internamente, R guarda un \codigo{Date} como el número de días
transcurridos desde el 1 de enero de 1970. Si escribes el data frame tal
cual, eso es lo que llega a la base. Hagamos la prueba en una base temporal
en memoria, leyendo el CSV de tiendas \emph{sin} convertir las fechas:

\begin{rcode}
# Base temporal para experimentar (no toca bi_retail.sqlite)
con_prueba <- dbConnect(RSQLite::SQLite(), ":memory:")

tiendas_r <- read_csv("datasets/tiendas.csv", show_col_types = FALSE)
class(tiendas_r$fecha_apertura)       # en R es una fecha

dbWriteTable(con_prueba, "tiendas", tiendas_r)
dbGetQuery(con_prueba, "
  SELECT nombre_tienda, fecha_apertura,
         typeof(fecha_apertura)             AS tipo,
         strftime('%Y', fecha_apertura)     AS anio
  FROM tiendas
  LIMIT 3
")
\end{rcode}

\begin{rsalida}
[1] "Date"
    nombre_tienda fecha_apertura tipo  anio
1 Sucursal Centro          14683 real -4672
2  Sucursal Norte          15511 real -4670
3    Sucursal Sur          15227 real -4671
\end{rsalida}

La fecha de apertura de la Sucursal Centro (15 de marzo de 2010) quedó
guardada como el número \codigo{14683} (días desde 1970), de tipo
\codigo{real}. Peor aún: las funciones de fecha de SQLite interpretan ese
número de otra forma (como un ``día juliano'', contado desde el año 4714
antes de Cristo), así que \codigo{strftime('\%Y', \ldots)} devuelve un año
absurdo, \codigo{-4672}, sin marcar ningún error. Es un error silencioso: el
más peligroso de todos.

\subsubsection{La solución: fechas como texto ISO}

La práctica recomendada es convertir las fechas a texto con formato ISO
\codigo{'AAAA-MM-DD'} antes de escribirlas, justo lo que hizo nuestra
función \codigo{leer\_para\_bd()}. Las funciones de fecha de SQLite
entienden ese formato de forma nativa:

\begin{rcode}
tiendas_ok <- tiendas_r %>%
  mutate(fecha_apertura = format(fecha_apertura, "%Y-%m-%d"))

dbWriteTable(con_prueba, "tiendas", tiendas_ok, overwrite = TRUE)
dbGetQuery(con_prueba, "
  SELECT nombre_tienda, fecha_apertura,
         typeof(fecha_apertura)             AS tipo,
         strftime('%Y', fecha_apertura)     AS anio
  FROM tiendas
  LIMIT 3
")

# De regreso en R, la columna es texto: la convertimos con as.Date()
dbReadTable(con_prueba, "tiendas") %>%
  mutate(fecha_apertura = as.Date(fecha_apertura)) %>%
  pull(fecha_apertura) %>%
  class()

dbDisconnect(con_prueba)
\end{rcode}

\begin{rsalida}
    nombre_tienda fecha_apertura tipo anio
1 Sucursal Centro     2010-03-15 text 2010
2  Sucursal Norte     2012-06-20 text 2012
3    Sucursal Sur     2011-09-10 text 2011
[1] "Date"
\end{rsalida}

Ahora la fecha se ve bien, es de tipo \codigo{text} y el año es correcto.
Al leerla de regreso en R llega como texto, y basta un \codigo{as.Date()}
(o \codigo{ymd()} de \paquete{lubridate}) para volver a tener una fecha.

\begin{consejo}
\paquete{RSQLite} ofrece una alternativa:
\codigo{dbConnect(RSQLite::SQLite(), ruta, extended\_types = TRUE)}. Con esa
opción, RSQLite guarda la fecha con el tipo declarado \codigo{DATE} y la
convierte de ida y vuelta a \codigo{Date} de R automáticamente. Es cómodo si
solo R va a leer la base, pero otras herramientas (Power BI, Python, un
visor de SQLite) seguirán viendo el número. El texto ISO funciona con
todas, por eso es la opción que usamos en el libro. El mismo cuidado aplica
a las horas (\codigo{hms}) y a las fechas con hora (\codigo{POSIXct}): se
guardan como segundos si no las conviertes a texto.
\end{consejo}

\subsubsection{Funciones de fecha de SQLite}

La función estrella es \codigo{strftime(formato, fecha)}, que extrae partes
de una fecha: \codigo{'\%Y'} año, \codigo{'\%m'} mes, \codigo{'\%d'} día,
\codigo{'\%Y-\%m'} año-mes, \codigo{'\%w'} día de la semana
(0 = domingo). También existen \codigo{date()} para sumar o restar
intervalos y \codigo{julianday()} para calcular diferencias en días.

\begin{ejemplo}{Ventas por mes}
La pregunta más común en BI: ¿cuánto vendimos cada mes?
\end{ejemplo}

\begin{rcode}
ventas_mes_sql <- dbGetQuery(con, "
  SELECT strftime('%Y-%m', fecha)         AS mes,
         COUNT(*)                         AS tickets,
         ROUND(SUM(total_transaccion), 2) AS ventas
  FROM transacciones
  GROUP BY mes
  ORDER BY mes
")
ventas_mes_sql
\end{rcode}

\begin{rsalida}
       mes tickets   ventas
1  2023-01      98 82559.66
2  2023-02      79 65899.34
3  2023-03      91 90930.61
4  2023-04      71 57037.74
5  2023-05      88 93910.21
6  2023-06      84 88509.12
7  2023-07      83 86456.05
8  2023-08      75 78929.49
9  2023-09      86 68855.87
10 2023-10      74 70004.10
11 2023-11      80 80506.89
12 2023-12      91 96393.61
\end{rsalida}

Diciembre es el mejor mes (\$96\,394) y abril el más flojo (\$57\,038, con
solo 71 tickets). Guardamos el resultado en \codigo{ventas\_mes\_sql} porque
más adelante lo compararemos con el cálculo equivalente en \paquete{dplyr}. Aunque la tabla ya trae una columna \codigo{mes}, es
buena práctica saber calcularla desde la fecha: muchas tablas reales solo
tienen la fecha.

\begin{ejemplo}{¿Qué día de la semana vendemos más?}
Operaciones quiere ajustar los turnos del personal según el día de la
semana.
\end{ejemplo}

\begin{rcode}
dbGetQuery(con, "
  SELECT strftime('%w', fecha) AS num_dia,   -- 0 = domingo
         CASE strftime('%w', fecha)
           WHEN '0' THEN 'domingo'   WHEN '1' THEN 'lunes'
           WHEN '2' THEN 'martes'    WHEN '3' THEN 'miércoles'
           WHEN '4' THEN 'jueves'    WHEN '5' THEN 'viernes'
           ELSE 'sábado'
         END AS dia,
         COUNT(*) AS tickets,
         ROUND(AVG(total_transaccion), 2) AS ticket_promedio
  FROM transacciones
  GROUP BY num_dia
  ORDER BY num_dia
")
\end{rcode}

\begin{rsalida}
  num_dia       dia tickets ticket_promedio
1       0   domingo     146         1018.51
2       1     lunes     136          947.63
3       2    martes     129          925.31
4       3 miércoles     155          911.80
5       4    jueves     146          995.07
6       5   viernes     158         1020.67
7       6    sábado     130          885.94
\end{rsalida}

Viernes y miércoles tienen más tickets (158 y 155); el viernes y el domingo
tienen además los tickets promedio más altos (alrededor de \$1\,020). El
sábado es el día más flojo en ticket promedio (\$886). Observa que \codigo{CASE} tiene una forma abreviada
(\codigo{CASE expresión WHEN valor THEN \ldots}) cuando comparas siempre la
misma expresión contra varios valores.

\begin{rcode}
# Antigüedad de los clientes al 31 de diciembre de 2023
dbGetQuery(con, "
  SELECT nombre, fecha_registro,
         CAST(julianday('2023-12-31') - julianday(fecha_registro)
              AS INTEGER)                    AS dias_antiguedad,
         date(fecha_registro, '+1 year')     AS primer_aniversario
  FROM clientes
  ORDER BY fecha_registro
  LIMIT 4
")
\end{rcode}

\begin{rsalida}
            nombre fecha_registro dias_antiguedad primer_aniversario
1  Carlos González     2020-01-03            1458         2021-01-03
2    Luis Martínez     2020-01-10            1451         2021-01-10
3 Patricia Ramírez     2020-01-15            1446         2021-01-15
4      Luis Flores     2020-01-17            1444         2021-01-17
\end{rsalida}

\codigo{julianday()} convierte la fecha a un número de días, así que la
resta da la diferencia en días; \codigo{CAST(\ldots{} AS INTEGER)} quita los
decimales. \codigo{date(fecha, '+1 year')} suma un año (también acepta
\codigo{'+7 days'}, \codigo{'-1 month'}, \codigo{'start of month'}, etc.).

\begin{hazlotu}
Calcula las ventas por trimestre usando la fecha. Pista: el mes se obtiene
con \codigo{CAST(strftime('\%m', fecha) AS INTEGER)} y el trimestre es
\codigo{(mes + 2) / 3} (en SQLite, la división entre enteros descarta los
decimales).
\end{hazlotu}

\subsection{JOIN: combinar tablas}

Casi ninguna pregunta de negocio se responde con una sola tabla. La tabla de
transacciones solo tiene \codigo{producto\_id}; para saber la
\emph{categoría} hay que ir a la tabla de productos. Eso es un
\termino{JOIN}. Los dos tipos que usarás el 95\% del tiempo son:

\begin{itemize}
  \item \codigo{INNER JOIN}: solo las filas que tienen pareja en ambas
        tablas (equivale a \codigo{inner\_join()}).
  \item \codigo{LEFT JOIN}: todas las filas de la tabla de la izquierda,
        tengan o no pareja; si no la tienen, las columnas de la derecha
        quedan en \codigo{NULL} (\codigo{NA} en R). Equivale a
        \codigo{left\_join()}.
\end{itemize}

\begin{figure}[H]
\centering
\begin{tikzpicture}[
  conj/.style={circle, minimum size=2.3cm, draw=primario, thick},
  etq/.style={font=\sffamily\small\bfseries, text=primario}
]
% INNER JOIN
\begin{scope}
  \begin{scope}
    \clip (0,0) circle (1.15cm);
    \fill[secundario!45] (1.3,0) circle (1.15cm);
  \end{scope}
  \node[conj] at (0,0) {};
  \node[conj] at (1.3,0) {};
  \node[font=\scriptsize\sffamily, text=gristexto] at (-0.45,0) {A};
  \node[font=\scriptsize\sffamily, text=gristexto] at (1.75,0) {B};
  \node[etq] at (0.65,-1.6) {INNER JOIN};
  \node[font=\scriptsize\sffamily, text=gristexto, align=center]
    at (0.65,-2.15) {solo filas con pareja\\en ambas tablas};
\end{scope}
% LEFT JOIN
\begin{scope}[xshift=6cm]
  \fill[secundario!45] (0,0) circle (1.15cm);
  \node[conj] at (0,0) {};
  \node[conj] at (1.3,0) {};
  \node[font=\scriptsize\sffamily, text=gristexto] at (-0.45,0) {A};
  \node[font=\scriptsize\sffamily, text=gristexto] at (1.75,0) {B};
  \node[etq] at (0.65,-1.6) {LEFT JOIN};
  \node[font=\scriptsize\sffamily, text=gristexto, align=center]
    at (0.65,-2.15) {todas las filas de A;\\NULL si no hay pareja en B};
\end{scope}
\end{tikzpicture}
\caption{Los dos tipos de JOIN más usados. A es la tabla después de
\texttt{FROM} y B la tabla después de \texttt{JOIN}.}
\label{fig:m5-joins}
\end{figure}

\begin{ejemplo}{Ventas por categoría de producto}
El gerente de compras quiere saber qué categorías venden más para
negociar mejor con los proveedores.
\end{ejemplo}

\begin{rcode}
dbGetQuery(con, "
  SELECT p.categoria,
         COUNT(*)                           AS tickets,
         SUM(t.cantidad)                    AS unidades,
         ROUND(SUM(t.total_transaccion), 0) AS ventas
  FROM transacciones AS t
  INNER JOIN productos AS p
          ON t.producto_id = p.producto_id  -- llave de unión
  GROUP BY p.categoria
  ORDER BY ventas DESC
")
\end{rcode}

\begin{rsalida}
     categoria tickets unidades ventas
1      Belleza     187      551 182143
2  Electrónica     137      403 132672
3     Juguetes     118      356 125093
4    Alimentos     115      356 121543
5   Automotriz     105      325 104085
6     Deportes      92      280  86171
7        Hogar      64      187  63407
8     Mascotas      74      201  57668
9       Libros      70      204  57162
10        Ropa      38      106  30048
\end{rsalida}

Fíjate en los alias de tabla: \codigo{transacciones AS t} permite escribir
\codigo{t.cantidad} en lugar de \codigo{transacciones.cantidad}. Cuando dos
tablas tienen una columna con el mismo nombre (aquí,
\codigo{producto\_id}), \emph{debes} indicar de qué tabla viene; si no, la
base responde con un error de columna ambigua. Belleza es por mucho la categoría líder
(\$182\,143, 187 tickets), seguida de Electrónica; Ropa queda al final con
\$30\,048. Parte de la explicación está en el catálogo: Belleza tiene 9
productos y Ropa solo 2.

\begin{ejemplo}{Clientes que nunca han comprado}
El CRM tiene 200 clientes registrados, pero solo 198 compraron. ¿Quiénes
son los que faltan? Son candidatos perfectos para una campaña de
activación.
\end{ejemplo}

\begin{rcode}
dbGetQuery(con, "
  SELECT c.cliente_id, c.nombre, c.ciudad, c.fecha_registro
  FROM clientes AS c
  LEFT JOIN transacciones AS t
         ON c.cliente_id = t.cliente_id
  WHERE t.transaccion_id IS NULL   -- sin pareja en transacciones
")
\end{rcode}

\begin{rsalida}
  cliente_id            nombre ciudad fecha_registro
1        171      Rosa Jiménez Mérida     2022-08-27
2        199 Ricardo Hernández Puebla     2023-09-16
\end{rsalida}

Este patrón (\codigo{LEFT JOIN} + \codigo{WHERE \ldots{} IS NULL}) es el
equivalente SQL de \codigo{anti\_join()}: encuentra las filas de la
izquierda que \emph{no} tienen pareja. Son Rosa Jiménez, de Mérida, y Ricardo
Hernández, de Puebla. Este último se registró en septiembre de 2023, así que
quizá solo necesita tiempo; Rosa lleva más de un año registrada sin
comprar.

Puedes encadenar varios JOIN en la misma consulta:

\begin{ejemplo}{Precio de venta contra precio de catálogo por zona}
Auditoría sospecha que en algunas zonas se vende muy por debajo del precio
de catálogo. Compara, por zona de la tienda, el precio promedio de venta
con el de catálogo.
\end{ejemplo}

\begin{rcode}
dbGetQuery(con, "
  SELECT ti.zona,
         COUNT(*)                         AS tickets,
         ROUND(AVG(t.precio_venta), 2)    AS precio_venta_prom,
         ROUND(AVG(p.precio_catalogo), 2) AS precio_catalogo_prom
  FROM transacciones AS t
  INNER JOIN productos AS p  ON t.producto_id = p.producto_id
  INNER JOIN tiendas   AS ti ON t.tienda_id   = ti.tienda_id
  GROUP BY ti.zona
  ORDER BY tickets DESC
")
\end{rcode}

\begin{rsalida}
    zona tickets precio_venta_prom precio_catalogo_prom
1  Norte     399            326.93               422.66
2 Centro     290            311.12               414.57
3  Oeste     109            306.93               394.45
4    Sur     108            303.23               425.21
5   Este      94            334.45               399.15
\end{rsalida}

En todas las zonas el precio promedio de venta (alrededor de \$300 a \$335)
está muy por debajo del precio promedio de catálogo (\$395 a \$425). Es un
\emph{hallazgo de calidad de datos} que ya conoces: en los datos del curso,
\codigo{precio\_venta} se generó sin relación con
\codigo{precio\_catalogo}. En una empresa real, una diferencia así
merecería una investigación con el área comercial antes de publicar
cualquier análisis de márgenes.

\begin{advertencia}
Un JOIN mal planteado puede \emph{multiplicar} filas. Si unes por una
columna que se repite en ambas tablas (por ejemplo, \codigo{ciudad} entre
clientes y tiendas), cada cliente se empareja con todas las tiendas de su
ciudad y tus sumas se inflan. Antes de sumar después de un JOIN, verifica
con \codigo{COUNT(*)} que el número de filas es el que esperas.
\end{advertencia}

\subsection{Subconsultas}

Una \termino{subconsulta} es una consulta dentro de otra, escrita entre
paréntesis. Sirve, por ejemplo, para comparar contra un valor que primero
hay que calcular.

\begin{ejemplo}{Productos más caros que el promedio}
Queremos la lista de productos cuyo precio de catálogo supera el promedio
del catálogo, con su margen.
\end{ejemplo}

\begin{rcode}
dbGetQuery(con, "
  SELECT nombre_producto, categoria, precio_catalogo, margen_pct
  FROM productos
  WHERE precio_catalogo > (SELECT AVG(precio_catalogo)
                           FROM productos)
  ORDER BY precio_catalogo DESC
  LIMIT 5
")

# ¿Cuál es ese promedio y cuántos productos lo superan?
dbGetQuery(con, "
  SELECT ROUND(AVG(precio_catalogo), 2) AS promedio,
         (SELECT COUNT(*) FROM productos
          WHERE precio_catalogo > (SELECT AVG(precio_catalogo)
                                   FROM productos)) AS arriba
  FROM productos
")
\end{rcode}

\begin{rsalida}
  nombre_producto   categoria precio_catalogo margen_pct
1     Producto 28    Mascotas          774.03      97.77
2     Producto 47  Automotriz          766.07      92.82
3     Producto 41 Electrónica          756.54      48.47
4     Producto 04    Juguetes          743.84      72.42
5     Producto 20     Belleza          741.63      81.42
  promedio arriba
1   416.92     25
\end{rsalida}

El precio promedio del catálogo es \$416.92 y 25 de los 50 productos lo
superan. Fíjate en la última consulta: una subconsulta puede ir incluso en
el \codigo{SELECT}, siempre que devuelva un solo valor.

Una subconsulta también puede ir en el \codigo{FROM}, como si fuera una
tabla temporal. Es útil para agregar en dos pasos:

\begin{rcode}
# Paso 1 (subconsulta): gasto total de cada cliente
# Paso 2: estadísticas sobre esos totales
dbGetQuery(con, "
  SELECT COUNT(*)                   AS clientes,
         ROUND(AVG(gasto), 2)       AS gasto_promedio,
         ROUND(MAX(gasto), 2)       AS gasto_maximo,
         SUM(gasto > 8000)          AS clientes_mas_de_8000
  FROM (SELECT cliente_id, SUM(total_transaccion) AS gasto
        FROM transacciones
        GROUP BY cliente_id) AS gasto_cliente
")
\end{rcode}

\begin{rsalida}
  clientes gasto_promedio gasto_maximo clientes_mas_de_8000
1      198        4848.45      14680.1                   23
\end{rsalida}

Los 198 clientes con compras gastaron en promedio \$4\,848 en el año; el
máximo es \$14\,680 (el cliente 5 que vimos antes) y 23 clientes superan los
\$8\,000. Cada subconsulta en el \codigo{FROM} necesita un alias (aquí,
\codigo{gasto\_cliente}).

\subsection{CTE: consultas legibles con WITH}

Cuando las subconsultas se anidan, el SQL se vuelve difícil de leer. Una
\termino{CTE} (\emph{Common Table Expression}) le pone nombre a un paso
intermedio con \codigo{WITH nombre AS (\ldots)} y luego lo usas como si
fuera una tabla. Es el equivalente SQL de guardar un resultado intermedio
en un objeto de R.

\begin{ejemplo}{Participación de cada tienda en las ventas}
¿Qué porcentaje de las ventas totales aporta cada sucursal?
\end{ejemplo}

\begin{rcode}
dbGetQuery(con, "
  WITH ventas_tienda AS (          -- paso 1: ventas por tienda
    SELECT tienda_id, SUM(total_transaccion) AS ventas
    FROM transacciones
    GROUP BY tienda_id
  ),
  total AS (                       -- paso 2: total de la empresa
    SELECT SUM(ventas) AS ventas_total FROM ventas_tienda
  )
  SELECT ti.nombre_tienda,         -- paso 3: combinar y calcular %
         ROUND(v.ventas, 0) AS ventas,
         ROUND(100.0 * v.ventas / total.ventas_total, 1) AS pct
  FROM ventas_tienda AS v
  INNER JOIN tiendas AS ti ON v.tienda_id = ti.tienda_id
  CROSS JOIN total
  ORDER BY v.ventas DESC
")
\end{rcode}

\begin{rsalida}
      nombre_tienda ventas  pct
1   Sucursal Outlet 114377 11.9
2      Sucursal Sur 105305 11.0
3  Sucursal Plaza A 105216 11.0
4  Sucursal Express 102791 10.7
5  Sucursal Premium 101538 10.6
6    Sucursal Norte  89024  9.3
7     Sucursal Este  88420  9.2
8   Sucursal Centro  88414  9.2
9    Sucursal Oeste  84672  8.8
10 Sucursal Plaza B  80236  8.4
\end{rsalida}

Las ventas están muy repartidas: la Sucursal Outlet lidera con 11.9\% y la
Plaza B cierra con 8.4\%. Llama la atención que el Outlet, la tienda más
pequeña después de Express (350~m²), sea la que más vende. \codigo{CROSS JOIN total} pega la única fila de
\codigo{total} a cada tienda. Y observa el \codigo{100.0}: en SQL, la
división entre dos enteros descarta los decimales (\codigo{5 / 2} da
\codigo{2}), así que multiplicar por \codigo{100.0} obliga a usar decimales.

\subsection{Funciones de ventana}
\label{sec:m5-ventanas}

Las \termino{funciones de ventana} calculan un valor para cada fila
mirando un grupo de filas relacionadas (la ``ventana''), \emph{sin colapsar}
las filas como hace \codigo{GROUP BY}. Son el equivalente de
\codigo{group\_by()} + \codigo{mutate()} en dplyr. Su sintaxis es:

\begin{sqlcode}
funcion() OVER (PARTITION BY grupo ORDER BY orden)
\end{sqlcode}

\codigo{PARTITION BY} define los grupos (como \codigo{group\_by()}) y
\codigo{ORDER BY} el orden dentro de cada grupo. Las más útiles en BI son
\codigo{ROW\_NUMBER()}, \codigo{RANK()}, \codigo{SUM() OVER} (acumulados) y
\codigo{LAG()} (valor de la fila anterior). SQLite las soporta desde la
versión 3.25; la que trae \paquete{RSQLite} es mucho más reciente.

\begin{ejemplo}{El producto estrella de cada categoría}
Para el catálogo impreso, Marketing quiere destacar el producto más vendido
(en dinero) de cada categoría.
\end{ejemplo}

\begin{rcode}
dbGetQuery(con, "
  WITH ventas_producto AS (
    SELECT p.categoria, p.nombre_producto,
           SUM(t.total_transaccion) AS ventas
    FROM transacciones AS t
    INNER JOIN productos AS p ON t.producto_id = p.producto_id
    GROUP BY p.categoria, p.nombre_producto
  ),
  ranking AS (
    SELECT categoria, nombre_producto, ROUND(ventas, 0) AS ventas,
           ROW_NUMBER() OVER (PARTITION BY categoria
                              ORDER BY ventas DESC) AS posicion
    FROM ventas_producto
  )
  SELECT * FROM ranking
  WHERE posicion = 1
  ORDER BY ventas DESC
")
\end{rcode}

\begin{rsalida}
     categoria nombre_producto ventas posicion
1     Juguetes     Producto 06  34028        1
2     Deportes     Producto 19  28368        1
3    Alimentos     Producto 31  25430        1
4   Automotriz     Producto 47  24527        1
5        Hogar     Producto 32  24333        1
6      Belleza     Producto 14  24023        1
7  Electrónica     Producto 21  22083        1
8       Libros     Producto 16  20501        1
9         Ropa     Producto 42  17448        1
10    Mascotas     Producto 22  16949        1
\end{rsalida}

\codigo{ROW\_NUMBER()} numera los productos de cada categoría de mayor a
menor venta, y el \codigo{WHERE} externo se queda con el número 1. No se
puede filtrar por una función de ventana en el mismo nivel donde se calcula;
por eso va en una CTE. Como la lista está ordenada por
ventas, su primera fila es también el producto más vendido de todo el
catálogo: el Producto 06, de Juguetes, con \$34\,028.

\begin{ejemplo}{Ventas acumuladas y crecimiento mensual}
La directora quiere ver, mes a mes, las ventas, el acumulado del año y el
crecimiento contra el mes anterior.
\end{ejemplo}

\begin{rcode}
dbGetQuery(con, "
  WITH mensual AS (
    SELECT strftime('%Y-%m', fecha) AS mes,
           SUM(total_transaccion)   AS ventas
    FROM transacciones
    GROUP BY mes
  )
  SELECT mes,
         ROUND(ventas, 0)                           AS ventas,
         ROUND(SUM(ventas) OVER (ORDER BY mes), 0)  AS acumulado,
         ROUND(LAG(ventas) OVER (ORDER BY mes), 0)  AS mes_anterior,
         ROUND(100.0 * (ventas - LAG(ventas) OVER (ORDER BY mes)) /
               LAG(ventas) OVER (ORDER BY mes), 1)  AS crec_pct
  FROM mensual
  ORDER BY mes
")
\end{rcode}

\begin{rsalida}
       mes ventas acumulado mes_anterior crec_pct
1  2023-01  82560     82560           NA       NA
2  2023-02  65899    148459        82560    -20.2
3  2023-03  90931    239390        65899     38.0
4  2023-04  57038    296427        90931    -37.3
5  2023-05  93910    390338        57038     64.6
6  2023-06  88509    478847        93910     -5.8
7  2023-07  86456    565303        88509     -2.3
8  2023-08  78929    644232        86456     -8.7
9  2023-09  68856    713088        78929    -12.8
10 2023-10  70004    783092        68856      1.7
11 2023-11  80507    863599        70004     15.0
12 2023-12  96394    959993        80507     19.7
\end{rsalida}

\codigo{SUM(ventas) OVER (ORDER BY mes)} suma todas las filas desde el
inicio hasta la actual: el acumulado de diciembre, \$959\,993, coincide con
el total del año. \codigo{LAG(ventas)} trae el valor del mes anterior (en
enero es \codigo{NA} porque no hay mes previo) y con él calculamos el
crecimiento: abril cayó 37.3\% y mayo se recuperó 64.6\%. Es exactamente lo
que en dplyr harías con \codigo{cumsum()} y \codigo{lag()}.

\begin{ejemplo}{Ranking de tiendas con empates}
¿Qué lugar ocupa cada tienda por número de tickets? Si dos tiendas empatan,
deben compartir lugar.
\end{ejemplo}

\begin{rcode}
dbGetQuery(con, "
  SELECT tienda_id,
         COUNT(*)                                   AS tickets,
         RANK()       OVER (ORDER BY COUNT(*) DESC) AS rank,
         DENSE_RANK() OVER (ORDER BY COUNT(*) DESC) AS dense_rank,
         ROW_NUMBER() OVER (ORDER BY COUNT(*) DESC) AS row_number
  FROM transacciones
  GROUP BY tienda_id
  ORDER BY tickets DESC
")
\end{rcode}

\begin{rsalida}
   tienda_id tickets rank dense_rank row_number
1          6     109    1          1          1
2          3     109    1          1          2
3         10     108    3          2          3
4          9     108    3          2          4
5          8     106    5          3          5
6          2      99    6          4          6
7          5      95    7          5          7
8          4      94    8          6          8
9          1      87    9          7          9
10         7      85   10          8         10
\end{rsalida}

Las tiendas 6 y 3 empatan con 109 tickets. \codigo{RANK()} les da el lugar
1 a ambas y salta al 3 (como en una competencia deportiva);
\codigo{DENSE\_RANK()} no deja huecos (1, 1, 2\ldots) y
\codigo{ROW\_NUMBER()} numera sin empates, desempatando de forma arbitraria.
Elige según la regla de negocio: para ``premiar a los tres primeros
lugares'' suele usarse \codigo{RANK()}; para ``quedarme exactamente con
una fila por grupo'', \codigo{ROW\_NUMBER()}.

\begin{hazlotu}
Modifica la consulta del producto estrella para mostrar los
\emph{tres} productos más vendidos de cada categoría (\codigo{posicion <= 3}).
¿Cuántas filas obtienes? ¿Por qué no son exactamente 30?
\end{hazlotu}

% ----------------------------------------------------------------------------
\section{Consultas parametrizadas e inyección SQL}
\label{sec:m5-parametros}

Muy pronto querrás que tus consultas dependan de un valor que cambia: la
ciudad que el usuario elige en un tablero de Shiny, el rango de fechas del
reporte de este mes, el cliente que se está consultando. La tentación es
armar el texto de la consulta con \codigo{paste()}. \textbf{No lo hagas}
con valores que vienen de fuera de tu código.

\subsection{El peligro de pegar texto: inyección SQL}

\begin{ejemplo}{Un buscador de clientes por ciudad}
Imagina un tablero donde el usuario escribe una ciudad en una caja de texto
y ve cuántos clientes hay. Un usuario malicioso (o simplemente descuidado)
escribe algo raro en la caja.
\end{ejemplo}

\begin{rcode}
# Lo que escribió el usuario en la caja de texto
ciudad_usuario <- "Mérida' OR '1'='1"

# MAL: pegar el texto del usuario dentro del SQL
sql_peligroso <- paste0(
  "SELECT COUNT(*) AS clientes FROM clientes WHERE ciudad = '",
  ciudad_usuario, "'"
)
cat(sql_peligroso, "\n")
dbGetQuery(con, sql_peligroso)
\end{rcode}

\begin{rsalida}
SELECT COUNT(*) AS clientes FROM clientes WHERE ciudad = 'Mérida' OR '1'='1' 
  clientes
1      200
\end{rsalida}

El usuario cerró la comilla de \codigo{'Mérida'} y agregó
\codigo{OR '1'='1'}, una condición que siempre es verdadera. Resultado: la
consulta devolvió los 200 clientes de la base en lugar de los de Mérida. A
esto se le llama \termino{inyección SQL}: el texto del usuario deja de ser
un \emph{dato} y se convierte en \emph{código} que la base ejecuta. Con un
poco más de malicia, el texto podría intentar algo destructivo:

\begin{sqlcode}
-- Si el usuario escribiera:   x'; DROP TABLE clientes; --
-- la consulta armada con paste() quedaría así:
SELECT COUNT(*) FROM clientes WHERE ciudad = 'x'; DROP TABLE clientes; --'
\end{sqlcode}

Según el motor y el driver, esa segunda instrucción podría ejecutarse y
borrar la tabla. La inyección SQL es una de las vulnerabilidades más
explotadas en la historia de la informática; no es un riesgo teórico.

\subsection{La solución: consultas parametrizadas}

En una \termino{consulta parametrizada} escribes el SQL con
\emph{marcadores} (\codigo{?}) donde irán los valores, y pasas los valores
por separado en el argumento \codigo{params}. La base recibe el SQL y los
datos por canales distintos, así que el valor \emph{nunca} se interpreta
como código, contenga lo que contenga.

\begin{rcode}
sql_seguro <- "SELECT COUNT(*) AS clientes FROM clientes WHERE ciudad = ?"

# Con el texto malicioso: busca literalmente esa "ciudad" y no la halla
dbGetQuery(con, sql_seguro, params = list(ciudad_usuario))

# Con un valor legítimo funciona normalmente
dbGetQuery(con, sql_seguro, params = list("Mérida"))
\end{rcode}

\begin{rsalida}
  clientes
1        0
  clientes
1       17
\end{rsalida}

Con el texto malicioso ahora obtenemos 0 clientes (no existe ninguna ciudad
que se llame así, con todo y comillas), y con \codigo{"Mérida"} obtenemos
la respuesta correcta. Además, los parámetros te ahorran pelearte con las
comillas: no tienes que escribir \codigo{'} alrededor de los textos.

También puedes usar marcadores con nombre (\codigo{:nombre}) y pasar una
lista con nombres, lo que hace más legibles las consultas con varios
parámetros:

\begin{ejemplo}{Reporte de ventas por método de pago en un periodo}
Finanzas quiere un reporte reutilizable: ventas de un método de pago entre
dos fechas cualesquiera.
\end{ejemplo}

\begin{rcode}
reporte_pago <- function(con, metodo, desde, hasta) {
  dbGetQuery(con, "
    SELECT metodo_pago,
           COUNT(*)                         AS tickets,
           ROUND(SUM(total_transaccion), 2) AS ventas
    FROM transacciones
    WHERE metodo_pago = :metodo
      AND fecha BETWEEN :desde AND :hasta
    GROUP BY metodo_pago",
    params = list(metodo = metodo,
                  desde  = format(as.Date(desde)),   # fecha como texto
                  hasta  = format(as.Date(hasta))))
}

reporte_pago(con, "Transferencia", "2023-01-01", "2023-03-31")
reporte_pago(con, "Efectivo", "2023-12-01", "2023-12-31")
\end{rcode}

\begin{rsalida}
    metodo_pago tickets   ventas
1 Transferencia      22 19523.79
  metodo_pago tickets   ventas
1    Efectivo      24 25212.12
\end{rsalida}

En el primer trimestre entraron 22 transferencias por \$19\,523.79, y en
diciembre hubo 24 pagos en efectivo por \$25\,212.12. La misma función sirve
para cualquier combinación sin tocar el SQL. Observa que convertimos las fechas con
\codigo{format(as.Date(\ldots))}: así la función acepta tanto texto como
objetos \codigo{Date} y siempre manda a la base el texto ISO que espera.

Si pasas en \codigo{params} un vector con varios valores, la consulta se
ejecuta una vez por valor y los resultados se apilan:

\begin{rcode}
dbGetQuery(con,
  "SELECT ciudad, COUNT(*) AS clientes, ROUND(AVG(edad), 1) AS edad_prom
   FROM clientes WHERE ciudad = ?",
  params = list(c("CDMX", "Puebla", "Mérida")))
\end{rcode}

\begin{rsalida}
  ciudad clientes edad_prom
1   CDMX       27      56.2
2 Puebla       28      51.3
3 Mérida       17      47.6
\end{rsalida}

\subsection{Alternativas: \texorpdfstring{\codigo{glue\_sql()}}{glue\_sql()} y \texorpdfstring{\codigo{sqlInterpolate()}}{sqlInterpolate()}}

A veces necesitas construir el texto SQL de forma dinámica (por ejemplo, una
lista \codigo{IN (\ldots)} de longitud variable, o el nombre de una tabla).
Para eso hay dos herramientas seguras que \emph{escapan} los valores (les
ponen las comillas correctas y neutralizan las comillas internas):

\begin{itemize}
  \item \codigo{glue::glue\_sql()}, del paquete \paquete{glue} (se instala
        con el tidyverse): usa llaves \codigo{\{\}} para insertar valores;
        \codigo{\{valor*\}} expande un vector separado por comas y
        \codigo{\{`nombre`\}} (con acentos graves) inserta un nombre de
        tabla o columna.
  \item \codigo{DBI::sqlInterpolate()}: usa marcadores \codigo{?nombre}.
\end{itemize}

\begin{rcode}
library(glue)

ciudades <- c("Monterrey", "Guadalajara", "Tijuana")
edad_min <- 40
sql_glue <- glue_sql("
  SELECT ciudad, COUNT(*) AS clientes
  FROM clientes
  WHERE ciudad IN ({ciudades*}) AND edad >= {edad_min}
  GROUP BY ciudad", .con = con)
sql_glue
dbGetQuery(con, sql_glue)

# El texto malicioso queda "desactivado": sus comillas se duplican
glue_sql("SELECT * FROM clientes WHERE ciudad = {ciudad_usuario}",
         .con = con)
sqlInterpolate(con, "SELECT * FROM clientes WHERE ciudad = ?ciudad",
               ciudad = ciudad_usuario)
\end{rcode}

\begin{rsalida}
<SQL> SELECT ciudad, COUNT(*) AS clientes
FROM clientes
WHERE ciudad IN ('Monterrey', 'Guadalajara', 'Tijuana') AND edad >= 40
GROUP BY ciudad
       ciudad clientes
1 Guadalajara       13
2   Monterrey       11
3     Tijuana        6
<SQL> SELECT * FROM clientes WHERE ciudad = 'Mérida'' OR ''1''=''1'
<SQL> SELECT * FROM clientes WHERE ciudad = 'Mérida'' OR ''1''=''1'
\end{rsalida}

\codigo{glue\_sql()} convirtió el vector \codigo{ciudades} en
\codigo{'Monterrey', 'Guadalajara', 'Tijuana'}, con sus comillas, y el número
40 sin comillas. Entre los clientes de 40 años o más, Guadalajara tiene 13,
Monterrey 11 y Tijuana 6. En las dos últimas consultas, las comillas simples del
texto malicioso se duplicaron (\codigo{''}), que en SQL es la forma de
escribir una comilla \emph{dentro} de un texto: todo el valor queda como un
solo texto inofensivo.

\begin{advertencia}
Regla de oro: \textbf{nunca} armes SQL con \codigo{paste()},
\codigo{paste0()} o \codigo{sprintf()} usando valores que escribe un usuario
o que vienen de un archivo externo. Usa \codigo{params}, y si necesitas
construir el texto, \codigo{glue\_sql()} o \codigo{sqlInterpolate()}. Esto
será especialmente importante en el Módulo~\ref{cap:shiny}, donde los
usuarios escriben en cajas de texto de tus tableros.
\end{advertencia}

\begin{hazlotu}
Escribe una función \codigo{clientes\_por\_segmento(con, segmento)} que
devuelva cuántos clientes premium y no premium hay en el segmento indicado
(\codigo{"Joven"}, \codigo{"Adulto"}, \codigo{"Maduro"} o
\codigo{"Senior"}), usando una consulta parametrizada.
\end{hazlotu}

% ----------------------------------------------------------------------------
\section{dbplyr: usar dplyr contra la base de datos}
\label{sec:m5-dbplyr}

¿Y si pudieras consultar la base con los mismos verbos de \paquete{dplyr}
que ya dominas, sin escribir SQL? Eso hace \paquete{dbplyr}: traduce tu
código de dplyr a SQL, lo manda a la base y solo te trae el resultado.

\subsection{Tablas remotas con \texorpdfstring{\codigo{tbl()}}{tbl()}}

\codigo{tbl(con, "tabla")} crea una \termino{tabla remota}: un objeto que
\emph{apunta} a la tabla dentro de la base pero no la descarga.

\begin{rcode}
transacciones_bd <- tbl(con, "transacciones")
productos_bd     <- tbl(con, "productos")

transacciones_bd %>%
  select(transaccion_id, fecha, producto_id, total_transaccion) %>%
  head(3)
\end{rcode}

%RECORTE 1-1,...,3-$
\begin{rsalida}
# Source:   SQL [3 x 4]
...
  transaccion_id fecha      producto_id total_transaccion
           <dbl> <chr>            <dbl>             <dbl>
1              1 2023-01-13           3             1216.
2              2 2023-01-09          47              232.
3              3 2023-02-20           6              826.
\end{rsalida}

La primera línea indica que el origen es una consulta SQL (de 3 filas y 4
columnas, por el \codigo{head(3)}); la segunda, que omitimos, indica el
motor y la ruta del archivo de la base. dbplyr no descargó la tabla: solo
pidió a la base las tres primeras filas para mostrártelas. Los importes se
ven redondeados (\codigo{1216.}) porque así imprime los números un tibble;
el valor completo sigue siendo 1215.88.

\subsection{Verbos de dplyr, \texorpdfstring{\codigo{show\_query()}}{show\_query()} y \texorpdfstring{\codigo{collect()}}{collect()}}

\begin{ejemplo}{Ventas por método de pago, versión dplyr}
Repitamos el reporte de métodos de pago, ahora con dplyr contra la base.
\end{ejemplo}

\begin{rcode}
consulta_pago <- transacciones_bd %>%
  group_by(metodo_pago) %>%
  summarise(tickets = n(),
            ventas  = round(sum(total_transaccion, na.rm = TRUE), 0)) %>%
  arrange(desc(ventas))

# ¿Qué SQL generó dbplyr?
show_query(consulta_pago)
\end{rcode}

\begin{rsalida}
<SQL>
SELECT
  `metodo_pago`,
  COUNT(*) AS `tickets`,
  ROUND(SUM(`total_transaccion`), 0) AS `ventas`
FROM `transacciones`
GROUP BY `metodo_pago`
ORDER BY `ventas` DESC
\end{rsalida}

\codigo{show\_query()} te enseña el SQL que dbplyr escribió por ti. Es una
excelente forma de \emph{aprender} SQL: escribe el código en dplyr y mira
cómo se traduce. Nota el \codigo{na.rm = TRUE}: en SQL las funciones de
agregación siempre ignoran los valores nulos, y dbplyr te lo recuerda con
un aviso si no lo escribes.

Para traer el resultado a R como un tibble normal usa \codigo{collect()}:

\begin{rcode}
resultado_pago <- collect(consulta_pago)
resultado_pago
class(resultado_pago)
\end{rcode}

\begin{rsalida}
# A tibble: 4 x 3
  metodo_pago     tickets ventas
  <chr>             <int>  <dbl>
1 Tarjeta Crédito     350 351171
2 Efectivo            315 303735
3 Tarjeta Débito      233 220855
4 Transferencia       102  84232
[1] "tbl_df"     "tbl"        "data.frame"
\end{rsalida}

Las cifras coinciden con las que obtuvimos con SQL escrito a mano. Ahora
\codigo{resultado\_pago} es un tibble en la memoria de R y puedes usarlo con
\paquete{ggplot2} o exportarlo a Excel.

\begin{concepto}{Evaluación diferida (\emph{lazy evaluation})}
Con dbplyr, cada verbo \emph{no ejecuta nada}: solo va construyendo la
receta de la consulta. La base trabaja únicamente cuando pides ver el
resultado (imprimirlo) o cuando llamas a \codigo{collect()}. Por eso puedes
encadenar muchos pasos sobre una tabla de 100 millones de filas sin
problema: todo el trabajo pesado lo hace el motor, y a R solo llega el
resultado final, normalmente pequeño.
\end{concepto}

Los \emph{joins} también se traducen. Repitamos las ventas por categoría:

\begin{rcode}
ventas_categoria_dbplyr <- transacciones_bd %>%
  inner_join(productos_bd, by = "producto_id") %>%
  group_by(categoria) %>%
  summarise(tickets = n(),
            ventas  = round(sum(total_transaccion, na.rm = TRUE), 0)) %>%
  arrange(desc(ventas)) %>%
  collect()

head(ventas_categoria_dbplyr, 4)
\end{rcode}

\begin{rsalida}
# A tibble: 4 x 3
  categoria   tickets ventas
  <chr>         <int>  <dbl>
1 Belleza         187 182143
2 Electrónica     137 132672
3 Juguetes        118 125093
4 Alimentos       115 121543
\end{rsalida}

Los números son idénticos a los del \codigo{INNER JOIN} en SQL: Belleza
lidera con \$182\,143.

\subsection{Límites de la traducción}

dbplyr traduce muchas funciones de R (aritmética, \codigo{if\_else()},
\codigo{case\_when()}, \codigo{substr()}, \codigo{paste0()},
\codigo{min\_rank()}, \codigo{lag()}, \codigo{cumsum()}, las funciones de
agregación\ldots), pero no todas. Algunas reglas prácticas:

\begin{rcode}
# Funciona: month() de lubridate se traduce a strftime() en SQLite
transacciones_bd %>%
  mutate(num_mes = month(fecha)) %>%
  count(num_mes) %>%
  head(3) %>%
  show_query()
\end{rcode}

\begin{rsalida}
<SQL>
SELECT `num_mes`, COUNT(*) AS `n`
FROM (
  SELECT
    `transacciones`.*,
    CAST(STRFTIME('%m', `fecha`) AS NUMERIC) AS `num_mes`
  FROM `transacciones`
) AS `q01`
GROUP BY `num_mes`
LIMIT 3
\end{rsalida}

\begin{itemize}
  \item Si una función no tiene traducción, dbplyr la pasa \emph{tal cual}
        a la base, y la base marca error de función inexistente. Por
        ejemplo, \codigo{floor\_date()} no se traduce en SQLite.
  \item Algunas funciones solo se traducen parcialmente. Por ejemplo,
        \codigo{str\_detect(x, "Tarjeta")} falla con
        \codigo{Only fixed patterns are supported on database backends};
        la solución es \codigo{str\_detect(x, fixed("Tarjeta"))}.
  \item Puedes escribir un fragmento de SQL literal con \codigo{sql()}:
        \codigo{mutate(mes = sql("strftime('\%Y-\%m', fecha)"))}.
  \item Si algo no se puede traducir, filtra y agrega primero en la base,
        haz \codigo{collect()} y termina el cálculo en R.
\end{itemize}

\subsection{¿SQL directo o dbplyr?}

Ambos caminos llegan al mismo resultado. La tabla~\ref{tab:m5-sql-dbplyr}
te ayuda a elegir.

\begin{table}[H]
\centering
\small
\caption{Cuándo conviene escribir SQL y cuándo usar dbplyr.}
\label{tab:m5-sql-dbplyr}
\begin{tabularx}{\textwidth}{@{}XX@{}}
\toprule
\textbf{Usa SQL directo (\codigo{dbGetQuery()}) cuando\ldots} &
\textbf{Usa dbplyr cuando\ldots} \\
\midrule
Ya tienes la consulta escrita (te la pasó el área de TI o viene de otro
reporte). & Quieres un solo estilo de código en todo tu proyecto. \\
Necesitas funciones propias del motor (\codigo{PRAGMA},
\codigo{EXPLAIN}, funciones de fecha especiales). & Estás explorando y
cambias la consulta a cada rato. \\
Creas o modificas estructura: \codigo{CREATE TABLE}, \codigo{CREATE VIEW},
\codigo{INSERT}, índices. & Quieres combinar pasos en la base con pasos en R
(\codigo{collect()} en medio). \\
La consulta la compartirás con personas que no usan R. & Prefieres que el
mismo código funcione con varios motores sin cambiar el dialecto. \\
\bottomrule
\end{tabularx}
\end{table}

\begin{consejo}
En la práctica, los analistas combinan ambos: dbplyr para explorar y
prototipar, y SQL escrito a mano para las consultas ``oficiales'' que se
guardan como vistas en el data warehouse. Si una consulta de dbplyr va a
producción, copia el resultado de \codigo{show\_query()}, revísalo y
guárdalo.
\end{consejo}

\begin{hazlotu}
Con dbplyr, calcula el ticket promedio por tienda (\codigo{tienda\_id}) solo
para pagos en efectivo, ordénalo de mayor a menor y mira el SQL con
\codigo{show\_query()} antes de hacer \codigo{collect()}.
\end{hazlotu}

% ----------------------------------------------------------------------------
\section{Diseño de un data warehouse en estrella}
\label{sec:m5-estrella}

Hasta ahora la base es una copia fiel de los CSV. Funciona, pero tiene
problemas típicos de los datos ``crudos'': no hay llaves declaradas, los
nombres de columnas son inconsistentes, faltan atributos útiles para
analizar (¿la fecha cae en fin de semana?, ¿cuánto costó lo vendido?) y
cualquiera podría insertar una venta de un producto inexistente. Un
\emph{data warehouse} resuelve esto con un diseño pensado para el análisis:
el \termino{esquema en estrella}.

\begin{concepto}{Hechos y dimensiones}
\begin{itemize}
  \item La \termino{tabla de hechos} (\emph{fact table}) registra los
        eventos del negocio que quieres medir, uno por fila, con sus
        \termino{métricas} numéricas: cantidad, importe, costo, margen.
        Tiene muchas filas y pocas columnas, casi todas numéricas o
        llaves.
  \item Las \termino{tablas de dimensión} describen el contexto de cada
        hecho: \emph{quién} compró (cliente), \emph{qué} (producto),
        \emph{dónde} (tienda), \emph{quién vendió} (vendedor) y
        \emph{cuándo} (fecha). Tienen pocas filas y muchas columnas de
        texto que sirven para filtrar y agrupar.
  \item La \termino{granularidad} (\emph{grain}) es el nivel de detalle de
        una fila de hechos. En nuestro caso: \emph{una fila por ticket de
        venta}.
\end{itemize}
Dibujado, la tabla de hechos queda al centro y las dimensiones alrededor,
como las puntas de una estrella (figura~\ref{fig:m5-estrella}).
\end{concepto}

\begin{figure}[H]
\centering
\begin{tikzpicture}[
  tabla/.style={draw=#1, rounded corners=2pt, fill=white,
    font=\ttfamily\scriptsize, align=left, inner sep=3pt,
    minimum width=3.2cm},
  titulo/.style={fill=#1, text=white, font=\sffamily\small\bfseries,
    minimum width=3.2cm, inner sep=3pt},
  rel/.style={-{Stealth[length=2.5mm]}, draw=gristexto, thick}
]
% Hechos
\node[titulo=acento] (ft) at (0,0) {fact\_ventas};
\node[tabla=acento, below=0pt of ft] (f) {\textbf{PK} venta\_id\\
\textbf{FK} fecha\_id\\ \textbf{FK} cliente\_id\\ \textbf{FK} producto\_id\\
\textbf{FK} tienda\_id\\ \textbf{FK} vendedor\_id\\ hora, metodo\_pago\\
cantidad, precio\_venta\\ total, costo\_total, margen};
% Dimensiones
\node[titulo=secundario] (dft) at (-5.3,1.3) {dim\_fecha};
\node[tabla=secundario, below=0pt of dft] (df) {\textbf{PK} fecha\_id\\
fecha, anio, trimestre\\ mes, nombre\_mes\\ dia\_semana, es\_fin\_semana};
\node[titulo=secundario] (dct) at (5.3,1.3) {dim\_cliente};
\node[tabla=secundario, below=0pt of dct] (dc) {\textbf{PK} cliente\_id\\
nombre, edad, genero\\ ciudad, segmento\\ es\_premium};
\node[titulo=secundario] (dpt) at (-5.3,-2.3) {dim\_producto};
\node[tabla=secundario, below=0pt of dpt] (dp) {\textbf{PK} producto\_id\\
nombre\_producto\\ categoria, marca\\ precio\_catalogo, costo};
\node[titulo=secundario] (dtt) at (5.3,-2.3) {dim\_tienda};
\node[tabla=secundario, below=0pt of dtt] (dt) {\textbf{PK} tienda\_id\\
nombre\_tienda, ciudad\\ zona, tipo, tamano\_m2};
\node[titulo=secundario] (dvt) at (0,-4.6) {dim\_vendedor};
\node[tabla=secundario, below=0pt of dvt] (dv) {\textbf{PK} vendedor\_id\\
nombre, puesto, sucursal};
% Relaciones
\draw[rel] ($(f.west)+(0,0.9)$) -- (df.east);
\draw[rel] ($(f.east)+(0,0.9)$) -- (dc.west);
\draw[rel] ($(f.west)+(0,-0.6)$) -- (dp.east);
\draw[rel] ($(f.east)+(0,-0.6)$) -- (dt.west);
\draw[rel] (f.south) -- (dvt.north);
\end{tikzpicture}
\caption{Esquema en estrella del data warehouse \texttt{bi\_retail}: la
tabla de hechos (naranja) al centro y cinco dimensiones alrededor.}
\label{fig:m5-estrella}
\end{figure}

\begin{negocio}{¿Por qué en estrella?}
El esquema en estrella es el estándar de facto en BI porque las consultas
son siempre iguales: ``suma las métricas de la tabla de hechos, filtrando y
agrupando por atributos de las dimensiones''. Power BI, Tableau y Qlik
están optimizados para él, los usuarios de negocio lo entienden de un
vistazo y los motores lo resuelven muy rápido. Si algún día trabajas con
Power BI, verás que su ``modelo de datos'' es justamente una estrella.
\end{negocio}

\subsection{Crear las tablas con llaves primarias y foráneas}

Esta vez no dejaremos que \codigo{dbWriteTable()} adivine los tipos:
escribiremos nosotros el \codigo{CREATE TABLE} de cada tabla con su llave
primaria, sus llaves foráneas y algunas reglas de validación
(\codigo{NOT NULL}, \codigo{CHECK}). A este código que define estructura se
le llama \termino{DDL} (\emph{Data Definition Language}). Lo guardamos en un
vector con nombre para poder reutilizarlo.

\begin{rcode}
ddl_estrella <- c(
  dim_fecha = "
    CREATE TABLE dim_fecha (
      fecha_id       INTEGER PRIMARY KEY,  -- AAAAMMDD, ej. 20230113
      fecha          TEXT NOT NULL UNIQUE, -- texto ISO 'AAAA-MM-DD'
      anio           INTEGER, trimestre INTEGER, mes INTEGER,
      nombre_mes     TEXT, dia INTEGER, dia_semana TEXT,
      num_dia_semana INTEGER,              -- 1 = lunes ... 7 = domingo
      es_fin_semana  INTEGER               -- 1 = sábado o domingo
    )",
  dim_producto = "
    CREATE TABLE dim_producto (
      producto_id INTEGER PRIMARY KEY, nombre_producto TEXT NOT NULL,
      categoria TEXT, subcategoria TEXT, marca TEXT,
      precio_catalogo REAL, costo REAL, proveedor TEXT, activo INTEGER
    )",
  dim_cliente = "
    CREATE TABLE dim_cliente (
      cliente_id INTEGER PRIMARY KEY, nombre TEXT NOT NULL,
      edad INTEGER, genero TEXT, ciudad TEXT, segmento TEXT,
      es_premium INTEGER, fecha_registro TEXT
    )",
  dim_tienda = "
    CREATE TABLE dim_tienda (
      tienda_id INTEGER PRIMARY KEY, nombre_tienda TEXT NOT NULL,
      ciudad TEXT, zona TEXT, tipo TEXT, tamano_m2 REAL,
      fecha_apertura TEXT
    )",
  dim_vendedor = "
    CREATE TABLE dim_vendedor (
      vendedor_id INTEGER PRIMARY KEY, nombre TEXT NOT NULL,
      departamento TEXT, puesto TEXT, sucursal TEXT, fecha_ingreso TEXT
    )",
  fact_ventas = "
    CREATE TABLE fact_ventas (
      venta_id     INTEGER PRIMARY KEY,
      fecha_id     INTEGER NOT NULL REFERENCES dim_fecha(fecha_id),
      cliente_id   INTEGER NOT NULL REFERENCES dim_cliente(cliente_id),
      producto_id  INTEGER NOT NULL REFERENCES dim_producto(producto_id),
      tienda_id    INTEGER NOT NULL REFERENCES dim_tienda(tienda_id),
      vendedor_id  INTEGER NOT NULL REFERENCES dim_vendedor(vendedor_id),
      hora TEXT, metodo_pago TEXT,
      cantidad     INTEGER NOT NULL CHECK (cantidad > 0),
      precio_venta REAL    NOT NULL CHECK (precio_venta > 0),
      total REAL NOT NULL, costo_total REAL, margen REAL
    )",
  log_etl = "
    CREATE TABLE log_etl (
      ejecucion_id  INTEGER PRIMARY KEY AUTOINCREMENT,
      fecha_hora    TEXT, archivo TEXT,
      filas_leidas  INTEGER, filas_validas INTEGER,
      filas_nuevas  INTEGER, estado TEXT, mensaje TEXT
    )"
)

crear_esquema <- function(con) {
  # SQLite trae apagada la revisión de llaves foráneas: la encendemos
  dbExecute(con, "PRAGMA foreign_keys = ON")
  for (sql in ddl_estrella) dbExecute(con, sql)
  invisible(names(ddl_estrella))
}

crear_esquema(con)
dbListTables(con)
\end{rcode}

\begin{rsalida}
 [1] "clientes"        "dim_cliente"     "dim_fecha"      
 [4] "dim_producto"    "dim_tienda"      "dim_vendedor"   
 [7] "empleados"       "fact_ventas"     "log_etl"        
[10] "productos"       "sqlite_sequence" "tiendas"        
[13] "transacciones"   "ventas_retail"  
\end{rsalida}

Ahora la base tiene las seis tablas crudas que cargamos al principio y las
nuevas tablas del \emph{data warehouse}, además de \codigo{log\_etl}, que
usaremos como bitácora del proceso de carga. La tabla
\codigo{sqlite\_sequence} la crea SQLite automáticamente para llevar la
cuenta del \codigo{AUTOINCREMENT}. Tres detalles del diseño:

\begin{itemize}
  \item \codigo{PRAGMA foreign\_keys = ON} es imprescindible en SQLite: sin
        él, las cláusulas \codigo{REFERENCES} se aceptan pero \emph{no se
        verifican}. Hay que activarlo en cada conexión nueva.
  \item La llave de la fecha es un entero con forma \codigo{AAAAMMDD}
        (\codigo{20230113}). Es una convención muy común en data warehouses:
        es compacta, se ordena bien y se lee a simple vista.
  \item Las columnas \codigo{CHECK (cantidad > 0)} hacen que la base
        rechace datos imposibles, sin importar quién o qué programa intente
        cargarlos.
\end{itemize}

\subsection{La dimensión de fechas, construida en R}

Una \termino{dimensión de fechas} (o calendario) tiene una fila por día con
todos los atributos que alguien podría querer para filtrar o agrupar: año,
trimestre, mes, nombre del mes, día de la semana, si es fin de semana, si es
feriado\ldots{} No viene de ningún sistema: la construimos nosotros, y
\paquete{lubridate} lo hace muy fácil.

\begin{rcode}
construir_dim_fecha <- function(desde, hasta) {
  tibble(fecha = seq(as.Date(desde), as.Date(hasta), by = "day")) %>%
    mutate(
      fecha_id   = as.integer(format(fecha, "%Y%m%d")),
      anio       = year(fecha),
      trimestre  = quarter(fecha),
      mes        = month(fecha),
      nombre_mes = as.character(month(fecha, label = TRUE, abbr = FALSE)),
      dia        = day(fecha),
      dia_semana = as.character(wday(fecha, label = TRUE, abbr = FALSE,
                                     week_start = 1)),
      num_dia_semana = wday(fecha, week_start = 1),   # 1 = lunes
      es_fin_semana  = as.integer(num_dia_semana >= 6),
      fecha      = format(fecha, "%Y-%m-%d")          # texto ISO
    ) %>%
    relocate(fecha_id)
}

dim_fecha <- construir_dim_fecha("2023-01-01", "2023-12-31")
nrow(dim_fecha)
dim_fecha %>% filter(fecha >= "2023-12-29") %>% select(-dia)
\end{rcode}

\begin{rsalida}
[1] 365
# A tibble: 3 x 9
  fecha_id fecha       anio trimestre   mes nombre_mes dia_semana
     <int> <chr>      <dbl>     <int> <dbl> <chr>      <chr>     
1 20231229 2023-12-29  2023         4    12 diciembre  viernes   
2 20231230 2023-12-30  2023         4    12 diciembre  sábado    
3 20231231 2023-12-31  2023         4    12 diciembre  domingo   
# i 2 more variables: num_dia_semana <dbl>, es_fin_semana <int>
\end{rsalida}

Tenemos 365 días con sus atributos. Los nombres de meses y días salen en
español porque R usa la configuración regional de tu sistema; si en tu
computadora salen en inglés, revisa \codigo{Sys.getlocale()}.

\subsection{Las demás dimensiones}

Las otras dimensiones salen de los CSV del curso con algunas
transformaciones: elegimos columnas, convertimos tipos (identificadores a
entero, lógicos a 0/1, fechas a texto ISO) y dejamos fuera los datos
personales que el análisis no necesita (el correo y el teléfono de los
clientes). Lo escribimos como una función que devuelve una lista con todas
las dimensiones listas para cargar:

\begin{rcode}
preparar_dimensiones <- function(carpeta = "datasets") {
  leer <- function(archivo) {
    read_csv(file.path(carpeta, archivo), show_col_types = FALSE)
  }
  empleados     <- leer("empleados.csv")
  transacciones <- leer("transacciones.csv")

  list(
    dim_fecha = construir_dim_fecha("2023-01-01", "2023-12-31"),

    dim_producto = leer("productos.csv") %>%
      transmute(producto_id = as.integer(producto_id), nombre_producto,
                categoria, subcategoria, marca, precio_catalogo, costo,
                proveedor, activo = as.integer(activo)),

    dim_cliente = leer("clientes.csv") %>%       # sin email ni teléfono
      transmute(cliente_id = as.integer(cliente_id), nombre,
                edad = as.integer(edad), genero, ciudad,
                segmento, es_premium = as.integer(es_premium),
                fecha_registro = format(fecha_registro, "%Y-%m-%d")),

    dim_tienda = leer("tiendas.csv") %>%
      transmute(tienda_id = as.integer(tienda_id), nombre_tienda, ciudad,
                zona, tipo, tamano_m2,
                fecha_apertura = format(fecha_apertura, "%Y-%m-%d")),

    # Los vendedores son los empleados que aparecen como vendedor_id
    dim_vendedor = empleados %>%
      semi_join(transacciones, by = c("empleado_id" = "vendedor_id")) %>%
      transmute(vendedor_id = as.integer(empleado_id), nombre,
                departamento, puesto, sucursal,
                fecha_ingreso = format(fecha_ingreso, "%Y-%m-%d"))
  )
}

dimensiones <- preparar_dimensiones()
map_int(dimensiones, nrow)
\end{rcode}

\begin{rsalida}
   dim_fecha dim_producto  dim_cliente   dim_tienda dim_vendedor 
         365           50          200           10           25 
\end{rsalida}

Ahora las cargamos. Aquí hay una diferencia importante con lo que hicimos al
principio: usamos \codigo{dbAppendTable()} y no \codigo{dbWriteTable()}.
\codigo{dbAppendTable()} \emph{agrega} filas a una tabla que ya existe,
respetando la estructura que diseñamos. En cambio,
\codigo{dbWriteTable(\ldots, overwrite = TRUE)} borraría la tabla y la
volvería a crear a su manera, sin llaves ni reglas de validación.

\begin{rcode}
cargar_dimensiones <- function(con, dimensiones) {
  dbBegin(con)                          # todo o nada (ver más adelante)
  for (tabla in names(dimensiones)) {
    dbAppendTable(con, tabla, dimensiones[[tabla]])
  }
  dbCommit(con)
  # Devolver cuántas filas quedaron en cada tabla
  map_int(names(dimensiones), function(tabla) {
    dbGetQuery(con, paste("SELECT COUNT(*) AS n FROM", tabla))$n
  }) %>% set_names(names(dimensiones))
}

cargar_dimensiones(con, dimensiones)
\end{rcode}

\begin{rsalida}
   dim_fecha dim_producto  dim_cliente   dim_tienda dim_vendedor 
         365           50          200           10           25 
\end{rsalida}

\begin{advertencia}
Un \emph{hallazgo de calidad de datos}: según el catálogo de empleados, de
los 25 vendedores solo unos cuantos pertenecen al departamento de Ventas;
los demás aparecen en Marketing, Finanzas, IT\ldots{} Compruébalo:
\end{advertencia}

\begin{rcode}
dbGetQuery(con, "
  SELECT departamento, COUNT(*) AS vendedores
  FROM dim_vendedor
  GROUP BY departamento
  ORDER BY vendedores DESC
")
\end{rcode}

\begin{rsalida}
         departamento vendedores
1 Atención al Cliente         10
2              Ventas          4
3                RRHH          4
4           Marketing          3
5            Finanzas          2
6         Operaciones          1
7                  IT          1
\end{rsalida}

Solo 4 de los 25 vendedores pertenecen al departamento de Ventas; 10 están
en Atención al Cliente y hay incluso uno de IT. Puede ser un error de los
datos simulados, o una realidad (personal de otras áreas que apoya en
caja). En una empresa lo reportarías a Recursos Humanos antes de publicar
un ranking de ``vendedores''. Lo importante es que un buen modelo de datos
hace visibles estas inconsistencias.

La base ya aplica las reglas que diseñamos. Si alguien intenta cargar un
producto con un identificador que ya existe, la llave primaria lo impide:

\begin{rcode}
resultado <- tryCatch(
  dbExecute(con, "INSERT INTO dim_producto (producto_id, nombre_producto)
                  VALUES (1, 'Producto repetido')"),
  error = function(e) conditionMessage(e)
)
resultado
\end{rcode}

\begin{rsalida}
[1] "UNIQUE constraint failed: dim_producto.producto_id"
\end{rsalida}

Usamos \codigo{tryCatch()} para \emph{atrapar} el error y mostrar su
mensaje sin detener el script: la base respondió que la restricción de
unicidad de \codigo{dim\_producto.producto\_id} falló. Es exactamente lo que
queremos.

\begin{hazlotu}
Intenta insertar en \codigo{dim\_cliente} una fila sin nombre
(\codigo{INSERT INTO dim\_cliente (cliente\_id) VALUES (999)}) dentro de un
\codigo{tryCatch()}. ¿Qué regla del diseño lo impide?
\end{hazlotu}

% ----------------------------------------------------------------------------
\section{ETL con R: extraer, transformar y cargar}
\label{sec:m5-etl}

Las dimensiones cambian poco; la tabla de hechos, en cambio, crece todos los
días. En una empresa, cada noche un proceso toma las ventas nuevas del
sistema de punto de venta y las agrega al \emph{data warehouse}. Ese
proceso es el \termino{ETL}:

\begin{figure}[H]
\centering
\begin{tikzpicture}[
  caja/.style={draw=#1, fill=#1!8, rounded corners=3pt, minimum width=3cm,
    minimum height=1.25cm, align=center, font=\sffamily\small},
  flecha/.style={-{Stealth[length=3mm]}, thick, draw=gristexto}
]
\node[caja=gristexto] (src) at (0,0) {\textbf{Fuentes}\\[-1pt]
  {\scriptsize CSV, Excel, sistemas}};
\node[caja=primario] (e) at (3.6,0) {\textbf{Extraer}\\[-1pt]
  {\scriptsize\ttfamily extraer\_*()}};
\node[caja=secundario] (t) at (7.2,0) {\textbf{Transformar}\\[-1pt]
  {\scriptsize\ttfamily transformar\_*()}};
\node[caja=acento] (l) at (10.8,0) {\textbf{Cargar}\\[-1pt]
  {\scriptsize\ttfamily cargar\_*()}};
\draw[flecha] (src) -- (e);
\draw[flecha] (e) -- (t);
\draw[flecha] (t) -- (l);
\node[font=\scriptsize\sffamily, text=gristexto, align=center, below=2pt of e]
  {leer sin modificar};
\node[font=\scriptsize\sffamily, text=gristexto, align=center, below=2pt of t]
  {limpiar, tipar,\\validar, llaves};
\node[font=\scriptsize\sffamily, text=gristexto, align=center, below=2pt of l]
  {transacción, sin\\duplicados, bitácora};
\end{tikzpicture}
\caption{Las tres etapas de un proceso ETL y las funciones de R que
escribiremos para cada una.}
\label{fig:m5-etl}
\end{figure}

Escribiremos cada etapa como una función pequeña y reutilizable, y una
función \codigo{ejecutar\_etl()} que las orquesta. Así, cuando llegue el
archivo de mañana, basta una línea para cargarlo.

\subsection{Transacciones: todo o nada}

Antes de escribir el ETL necesitas un concepto clave. Una
\termino{transacción} agrupa varias operaciones para que se apliquen
\emph{todas o ninguna}. Piensa en un pedido: se guarda el encabezado (quién
compra) y después los renglones (qué compra). Si los renglones fallan, no
quieres un encabezado huérfano en la base.

\begin{ejemplo}{Un pedido que falla a la mitad}
Probemos en una base en memoria qué pasa sin transacción y con transacción.
\end{ejemplo}

\begin{rcode}
con_pedidos <- dbConnect(RSQLite::SQLite(), ":memory:")
dbExecute(con_pedidos, "PRAGMA foreign_keys = ON")
dbExecute(con_pedidos, "CREATE TABLE productos (id INTEGER PRIMARY KEY)")
dbExecute(con_pedidos, "INSERT INTO productos VALUES (1), (2)")
dbExecute(con_pedidos, "CREATE TABLE pedidos (pedido_id INTEGER
                          PRIMARY KEY, cliente TEXT)")
dbExecute(con_pedidos, "CREATE TABLE renglones (pedido_id INTEGER,
            producto_id INTEGER REFERENCES productos(id))")

guardar_pedido <- function(con, id, cliente, productos) {
  dbAppendTable(con, "pedidos", tibble(pedido_id = id,
                                       cliente = cliente))
  dbAppendTable(con, "renglones", tibble(pedido_id = id,
                                         producto_id = productos))
}

# 1) SIN transacción: el producto 99 no existe y los renglones fallan
try(guardar_pedido(con_pedidos, 1, "Ana", c(1, 99)), silent = TRUE)
dbGetQuery(con_pedidos, "SELECT * FROM pedidos")   # ¡encabezado huérfano!
\end{rcode}

\begin{rsalida}
[1] 0
[1] 0
[1] 2
[1] 0
[1] 0
  pedido_id cliente
1         1     Ana
\end{rsalida}

El encabezado del pedido 1 se guardó, pero sus renglones no: la base quedó
inconsistente. Ahora con transacción. \codigo{dbBegin()} la inicia,
\codigo{dbCommit()} confirma los cambios y \codigo{dbRollback()} los
deshace todos:

\begin{rcode}
# 2) CON transacción: si algo falla, se deshace todo
dbBegin(con_pedidos)
tryCatch({
  guardar_pedido(con_pedidos, 2, "Luis", c(1, 99))
  dbCommit(con_pedidos)
}, error = function(e) {
  dbRollback(con_pedidos)
  cat("Se canceló el pedido:", conditionMessage(e), "\n")
})
dbGetQuery(con_pedidos, "SELECT * FROM pedidos")   # el 2 no quedó
dbDisconnect(con_pedidos)
\end{rcode}

\begin{rsalida}
Se canceló el pedido: FOREIGN KEY constraint failed 
  pedido_id cliente
1         1     Ana
\end{rsalida}

Con la transacción, el pedido 2 no dejó rastro: la base quedó exactamente
como estaba antes. \paquete{DBI} también ofrece
\codigo{dbWithTransaction(con, \{ \ldots{} \})}, que hace el
\codigo{dbBegin()}, \codigo{dbCommit()} y \codigo{dbRollback()} por ti.

\subsection{Extraer}

La etapa de extracción lee los datos \emph{sin modificarlos}. En nuestro
caso, archivos CSV que el sistema de punto de venta exporta cada cierto
tiempo. Forzamos la hora a texto con \codigo{col\_types} para conservarla
tal como viene (\codigo{"13:15"}).

\begin{rcode}
extraer_transacciones <- function(ruta) {
  read_csv(ruta, show_col_types = FALSE,
           col_types = cols(hora = col_character()))
}

# Otra fuente frecuente: una hoja de Excel
extraer_excel <- function(ruta, hoja) {
  readxl::read_excel(ruta, sheet = hoja)
}
\end{rcode}

Para simular la operación real, partimos \archivo{transacciones.csv} en dos
``exportaciones'' del sistema: el primer semestre y otra que va de mayo a
diciembre. A propósito, las dos se \emph{traslapan} en mayo y junio, algo
muy común cuando un sistema reenvía datos por seguridad. Nuestro ETL debe
evitar cargar esas ventas dos veces.

\begin{rcode}
dir.create("datos/entrada", showWarnings = FALSE)
trans_origen <- extraer_transacciones("datasets/transacciones.csv")

trans_origen %>% filter(fecha < as.Date("2023-07-01")) %>%
  write_csv("datos/entrada/ventas_2023_s1.csv")
trans_origen %>% filter(fecha >= as.Date("2023-05-01")) %>%
  write_csv("datos/entrada/ventas_2023_s2.csv")

list.files("datos/entrada")
\end{rcode}

\begin{rsalida}
[1] "ventas_2023_12_31.csv" "ventas_2023_s1.csv"   
[3] "ventas_2023_s2.csv"   
\end{rsalida}

\subsection{Transformar}

La transformación convierte los datos crudos en filas válidas para
\codigo{fact\_ventas}: descarta registros imposibles o duplicados, calcula
la llave de fecha, convierte tipos, agrega el costo desde la dimensión de
productos y calcula las métricas. Nota que \emph{recalculamos} el total
(\codigo{cantidad * precio\_venta}) en lugar de confiar en la columna del
archivo: en un ETL las métricas derivadas se calculan con una regla única.

\begin{rcode}
transformar_transacciones <- function(datos, costos) {
  datos %>%
    # 1. Calidad: quitar filas sin llave, imposibles o repetidas
    filter(!is.na(transaccion_id), cantidad > 0, precio_venta > 0) %>%
    distinct(transaccion_id, .keep_all = TRUE) %>%
    # 2. Enriquecer con el costo unitario del producto
    left_join(costos, by = "producto_id") %>%
    # 3. Tipos, llaves y métricas con los nombres del data warehouse
    transmute(
      venta_id    = as.integer(transaccion_id),
      fecha_id    = as.integer(format(as.Date(fecha), "%Y%m%d")),
      cliente_id  = as.integer(cliente_id),
      producto_id = as.integer(producto_id),
      tienda_id   = as.integer(tienda_id),
      vendedor_id = as.integer(vendedor_id),
      hora, metodo_pago,
      cantidad    = as.integer(cantidad),
      precio_venta,
      total       = round(cantidad * precio_venta, 2),
      costo_total = round(cantidad * costo, 2),
      margen      = round(total - costo_total, 2)
    )
}

# Prueba rápida con las primeras filas del archivo
costos <- dbGetQuery(con, "SELECT producto_id, costo FROM dim_producto")
trans_origen %>% head(3) %>% transformar_transacciones(costos) %>%
  select(venta_id, fecha_id, cantidad, total, costo_total, margen)
\end{rcode}

\begin{rsalida}
# A tibble: 3 x 6
  venta_id fecha_id cantidad total costo_total margen
     <int>    <int>    <int> <dbl>       <dbl>  <dbl>
1        1 20230113        4 1216.       1177.   38.7
2        2 20230109        4  232.        220    12.3
3        3 20230220        3  826.        254.  573. 
\end{rsalida}

La venta 1 (4 piezas a \$303.97) suma \$1\,215.88, le costó a la empresa
\$1\,177 y dejó un margen de solo \$38.7; la venta 3, en cambio, dejó
\$573 de margen sobre \$826. El tibble redondea al imprimir, pero en la
tabla se guardan los valores con dos decimales. Probar la transformación
con unas cuantas filas antes de cargar es una buena costumbre.

\subsection{Cargar sin duplicados}

Una \termino{carga incremental} agrega solo lo nuevo. Hay dos formas
habituales de evitar duplicados:

\begin{enumerate}
  \item \textbf{En R}, con \codigo{anti\_join()}: consultas qué llaves ya
        existen en la base y te quedas con las filas cuya llave no está.
  \item \textbf{En la base}, con \codigo{INSERT OR IGNORE}: como
        \codigo{venta\_id} es llave primaria, SQLite ignora en silencio las
        filas cuya llave ya existe. (En PostgreSQL el equivalente es
        \codigo{INSERT \ldots{} ON CONFLICT DO NOTHING}; en SQL Server, un
        \codigo{MERGE} o un \codigo{WHERE NOT EXISTS}.)
\end{enumerate}

\begin{rcode}
# Opción 1: evitar duplicados en R con anti_join()
cargar_hechos <- function(con, datos) {
  existentes <- dbGetQuery(con, "SELECT venta_id FROM fact_ventas")
  nuevos <- anti_join(datos, existentes, by = "venta_id")
  if (nrow(nuevos) > 0) dbAppendTable(con, "fact_ventas", nuevos)
  nrow(nuevos)                  # devuelve cuántas filas se cargaron
}

# Opción 2: dejar que la base ignore las llaves repetidas
cargar_hechos_sql <- function(con, datos) {
  dbWriteTable(con, "stg_ventas", datos, temporary = TRUE,
               overwrite = TRUE)          # tabla temporal de paso
  n <- dbExecute(con, "INSERT OR IGNORE INTO fact_ventas
                       SELECT * FROM stg_ventas")
  dbExecute(con, "DROP TABLE stg_ventas")
  n
}
\end{rcode}

La opción 1 es muy clara y funciona bien con miles o cientos de miles de
filas. La opción 2 escala mejor (no trae las llaves a R) y usa una
\termino{tabla de paso} o \emph{staging}: una tabla temporal donde se
depositan los datos antes de pasarlos a la tabla definitiva. La tabla
temporal solo existe mientras dura la conexión.

\begin{advertencia}
En SQLite, \codigo{INSERT OR IGNORE} ignora en silencio \emph{cualquier}
fila que viole una restricción, no solo las de llave repetida: también las
que no cumplen un \codigo{CHECK} o un \codigo{NOT NULL}. Una venta con
cantidad negativa simplemente ``desaparecería'' sin error. Si quieres que
esas filas detengan la carga, usa la opción 1 o un \codigo{INSERT \ldots{}
SELECT \ldots{} WHERE venta\_id NOT IN (SELECT venta\_id FROM
fact\_ventas)}. En nuestro ETL, además, las reglas de calidad de
\codigo{transformar\_transacciones()} ya descartan esas filas antes de
llegar a la base.
\end{advertencia}

\subsection{Bitácora y orquestación: \texorpdfstring{\codigo{ejecutar\_etl()}}{ejecutar\_etl()}}

Un ETL profesional deja rastro de cada ejecución: cuándo corrió, qué archivo
procesó, cuántas filas leyó, cuántas cargó y si terminó bien o con error.
Esa \termino{bitácora} (\emph{log}) es lo primero que revisas cuando
alguien pregunta ``¿por qué el tablero no tiene las ventas de ayer?''.

\begin{rcode}
registrar_log <- function(con, archivo, leidas, validas, nuevas,
                          estado, mensaje = NA_character_) {
  dbAppendTable(con, "log_etl", tibble(
    fecha_hora    = format(Sys.time(), "%Y-%m-%d %H:%M:%S"),
    archivo       = basename(archivo),
    filas_leidas  = leidas,  filas_validas = validas,
    filas_nuevas  = nuevas,  estado = estado,  mensaje = mensaje
  ))
}

ejecutar_etl <- function(con, ruta, cargar = cargar_hechos) {
  leidas <- validas <- NA_integer_
  en_transaccion <- FALSE
  estado <- tryCatch({
    datos   <- extraer_transacciones(ruta)                  # E
    leidas  <- nrow(datos)
    costos  <- dbGetQuery(con, "SELECT producto_id, costo
                                FROM dim_producto")
    limpios <- transformar_transacciones(datos, costos)     # T
    validas <- nrow(limpios)
    dbBegin(con); en_transaccion <- TRUE
    nuevas  <- cargar(con, limpios)                         # L
    registrar_log(con, ruta, leidas, validas, nuevas, "OK")
    dbCommit(con)
    sprintf("OK: %d leídas, %d nuevas", leidas, nuevas)
  }, error = function(e) {
    if (en_transaccion) dbRollback(con)     # deshacer la carga a medias
    registrar_log(con, ruta, leidas, validas, 0L, "ERROR",
                  conditionMessage(e))
    paste("ERROR:", conditionMessage(e))
  })
  cat(basename(ruta), "->", estado, "\n")
  invisible(estado)
}
\end{rcode}

Observa cómo encajan las piezas: si cualquier paso falla, el bloque
\codigo{error} deshace la transacción (si ya había empezado) y registra el
error en la bitácora; la función nunca deja la base a medias. Ahora
ejecutemos el ETL como lo haría el proceso nocturno:

\begin{rcode}
ejecutar_etl(con, "datos/entrada/ventas_2023_s1.csv")
ejecutar_etl(con, "datos/entrada/ventas_2023_s2.csv")

# ¿Y si por error se vuelve a ejecutar el mismo archivo?
ejecutar_etl(con, "datos/entrada/ventas_2023_s2.csv")

# Lo mismo con la opción 2 (INSERT OR IGNORE)
ejecutar_etl(con, "datos/entrada/ventas_2023_s2.csv",
             cargar = cargar_hechos_sql)

dbGetQuery(con, "SELECT COUNT(*) AS filas FROM fact_ventas")
\end{rcode}

\begin{rsalida}
ventas_2023_s1.csv -> OK: 511 leídas, 511 nuevas 
ventas_2023_s2.csv -> OK: 661 leídas, 489 nuevas 
ventas_2023_s2.csv -> OK: 661 leídas, 0 nuevas 
ventas_2023_s2.csv -> OK: 661 leídas, 0 nuevas 
  filas
1  1000
\end{rsalida}

El primer archivo cargó sus 511 ventas. El segundo leyó 661, pero solo 489
eran nuevas: las 172 de mayo y junio ya estaban y se descartaron. Las dos
ejecuciones repetidas (una con \codigo{anti\_join()} y otra con
\codigo{INSERT OR IGNORE}) no cargaron nada. Resultado: exactamente 1\,000
filas, las mismas del archivo original, sin un solo duplicado.

Probemos ahora el caso de error. Llega un archivo con tres ventas nuevas,
pero una tiene un producto que no existe en el catálogo (el 999):

\begin{rcode}
lote_malo <- trans_origen %>%
  slice(1:3) %>%
  mutate(transaccion_id = 1001:1003, fecha = as.Date("2023-12-31"),
         producto_id = c(5, 999, 7))
write_csv(lote_malo, "datos/entrada/ventas_2023_12_31.csv")

ejecutar_etl(con, "datos/entrada/ventas_2023_12_31.csv")
dbGetQuery(con, "SELECT COUNT(*) AS filas FROM fact_ventas")
\end{rcode}

\begin{rsalida}
ventas_2023_12_31.csv -> ERROR: FOREIGN KEY constraint failed 
  filas
1  1000
\end{rsalida}

La llave foránea hacia \codigo{dim\_producto} rechazó el lote, la
transacción se deshizo y la tabla de hechos sigue con 1\,000 filas: ni
siquiera las dos ventas válidas del archivo entraron. ¿Es eso lo que
queremos? En la mayoría de los data warehouses, sí: es preferible cargar el
archivo completo una vez corregido que tener un día ``a medias''. Veamos la
bitácora (omitimos la columna de fecha y hora, que cambia en cada
ejecución):

\begin{rcode}
dbGetQuery(con, "
  SELECT ejecucion_id AS id, archivo, filas_leidas AS leidas,
         filas_validas AS validas, filas_nuevas AS nuevas, estado
  FROM log_etl
")
dbGetQuery(con, "SELECT mensaje FROM log_etl WHERE estado = 'ERROR'")
\end{rcode}

\begin{rsalida}
  id               archivo leidas validas nuevas estado
1  1    ventas_2023_s1.csv    511     511    511     OK
2  2    ventas_2023_s2.csv    661     661    489     OK
3  3    ventas_2023_s2.csv    661     661      0     OK
4  4    ventas_2023_s2.csv    661     661      0     OK
5  5 ventas_2023_12_31.csv      3       3      0  ERROR
                        mensaje
1 FOREIGN KEY constraint failed
\end{rsalida}

La bitácora cuenta la historia completa de las cargas. En un proceso
programado (por ejemplo, con el Programador de tareas de Windows o
\codigo{cron}, que verás en el Módulo~\ref{cap:reportes}), revisarías esta
tabla cada mañana o, mejor, harías que el propio script te avise por
correo si hay un \codigo{ERROR}.

\begin{consejo}
Tres reglas de oro de un ETL confiable: \textbf{(1)} que sea
\emph{idempotente}, es decir, que ejecutarlo dos veces con el mismo archivo
deje la base igual que ejecutarlo una vez (lo logramos con la llave primaria
y el \codigo{anti\_join()}); \textbf{(2)} que sea \emph{atómico}: o carga
todo el lote o nada (la transacción); y \textbf{(3)} que sea
\emph{trazable}: cada ejecución deja registro (la bitácora).
\end{consejo}

\subsection{Extraer desde Excel}

Muchas fuentes llegan en Excel. El archivo \archivo{datos\_empresa.xlsx} del
curso tiene una hoja \codigo{Tiendas}. ¿Coincide con lo que cargamos en
\codigo{dim\_tienda}? Es una validación típica de ETL:

\begin{rcode}
tiendas_excel <- extraer_excel("datasets/datos_empresa.xlsx", "Tiendas")
class(tiendas_excel$fecha_apertura)

tiendas_excel_t <- tiendas_excel %>%
  transmute(tienda_id = as.integer(tienda_id), nombre_tienda,
            fecha_apertura = format(as.Date(fecha_apertura), "%Y-%m-%d"))

# Filas del Excel que no coinciden exactamente con la dimensión
tiendas_excel_t %>%
  anti_join(dbReadTable(con, "dim_tienda"),
            by = c("tienda_id", "nombre_tienda", "fecha_apertura"))
\end{rcode}

\begin{rsalida}
[1] "POSIXct" "POSIXt" 
# A tibble: 0 x 3
# i 3 variables: tienda_id <int>, nombre_tienda <chr>,
#   fecha_apertura <chr>
\end{rsalida}

\codigo{readxl} lee las fechas de Excel como fecha con hora
(\codigo{POSIXct}); por eso las pasamos a \codigo{Date} y luego a texto ISO
antes de comparar. El \codigo{anti\_join()} no devolvió ninguna fila: la
hoja de Excel y la dimensión coinciden.

\begin{hazlotu}
Crea un archivo \archivo{datos/entrada/ventas\_2023\_12\_31\_v2.csv}
corrigiendo el lote malo (cambia el producto 999 por el 9) y ejecútalo con
\codigo{ejecutar\_etl()}. ¿Cuántas filas nuevas se cargan? ¿Qué muestra la
bitácora? (Si lo haces, recuerda que la tabla de hechos tendrá 1\,003 filas
y los resultados de las siguientes secciones cambiarán ligeramente.)
\end{hazlotu}

% ----------------------------------------------------------------------------
\section{Vistas e índices: el data warehouse listo para BI}
\label{sec:m5-vistas}

\subsection{Vistas para los KPIs}

Una \termino{vista} (\codigo{CREATE VIEW}) es una consulta guardada en la
base con un nombre; se consulta como si fuera una tabla, pero no guarda
datos: cada vez que la consultas, se ejecuta la consulta con los datos
actuales. Son ideales para publicar los KPIs ``oficiales'': todos (tus
scripts de R, Power BI, el tablero de Shiny) leen la misma definición.

\begin{rcode}
vistas <- c(
  # Vista "ancha": cada venta con todos sus atributos descriptivos
  "CREATE VIEW vw_ventas_detalle AS
   SELECT f.*, d.fecha, d.anio, d.mes, d.nombre_mes, d.dia_semana,
          d.es_fin_semana, p.nombre_producto, p.categoria, p.marca,
          c.nombre AS cliente, c.ciudad AS ciudad_cliente, c.segmento,
          c.es_premium, t.nombre_tienda, t.ciudad AS ciudad_tienda,
          t.zona, t.tipo AS tipo_tienda, v.nombre AS vendedor
   FROM fact_ventas AS f
   JOIN dim_fecha    AS d ON f.fecha_id    = d.fecha_id
   JOIN dim_producto AS p ON f.producto_id = p.producto_id
   JOIN dim_cliente  AS c ON f.cliente_id  = c.cliente_id
   JOIN dim_tienda   AS t ON f.tienda_id   = t.tienda_id
   JOIN dim_vendedor AS v ON f.vendedor_id = v.vendedor_id",
  # KPIs mensuales
  "CREATE VIEW vw_kpi_mensual AS
   SELECT d.anio, d.mes, d.nombre_mes,
          COUNT(*)                     AS tickets,
          COUNT(DISTINCT f.cliente_id) AS clientes,
          ROUND(SUM(f.total), 2)       AS ventas,
          ROUND(SUM(f.margen), 2)      AS margen,
          ROUND(AVG(f.total), 2)       AS ticket_promedio
   FROM fact_ventas AS f
   JOIN dim_fecha AS d ON f.fecha_id = d.fecha_id
   GROUP BY d.anio, d.mes, d.nombre_mes"
)
for (sql in vistas) dbExecute(con, sql)

dbGetQuery(con, "SELECT * FROM vw_kpi_mensual") %>% head(6)
\end{rcode}

\begin{rsalida}
  anio mes nombre_mes tickets clientes   ventas   margen
1 2023   1      enero      98       80 82559.66 34415.12
2 2023   2    febrero      79       67 65899.34 29212.36
3 2023   3      marzo      91       72 90930.61 46572.20
4 2023   4      abril      71       57 57037.74 19817.86
5 2023   5       mayo      88       73 93910.21 49550.60
6 2023   6      junio      84       68 88509.12 42049.19
  ticket_promedio
1          842.45
2          834.17
3          999.24
4          803.35
5         1067.16
6         1053.68
\end{rsalida}

La vista entrega, mes por mes, tickets, clientes distintos, ventas, margen
y ticket promedio. Marzo y mayo destacan por su margen (\$46\,572 y
\$49\,551), mientras que abril es el mes más débil también en margen
(\$19\,818). La vista ancha \codigo{vw\_ventas\_detalle} es
útil para analistas que no quieren escribir \emph{joins}: consultan una sola
``tabla'' con todo.

\begin{rcode}
dbGetQuery(con, "
  SELECT zona, ROUND(SUM(total), 0) AS ventas,
         ROUND(SUM(margen), 0) AS margen,
         ROUND(100.0 * SUM(margen) / SUM(total), 1) AS margen_pct
  FROM vw_ventas_detalle
  GROUP BY zona
  ORDER BY ventas DESC
")
\end{rcode}

\begin{rsalida}
    zona ventas margen margen_pct
1  Norte 388941 174227       44.8
2 Centro 275878 126667       45.9
3  Oeste 105216  48086       45.7
4    Sur 101538  49097       48.4
5   Este  88420  46559       52.7
\end{rsalida}

La zona Norte concentra las ventas (\$388\,941, con cuatro tiendas), pero su
margen porcentual (44.8\%) es el más bajo; la zona Este, con una sola
tienda, tiene el margen más alto (52.7\%). Con una sola consulta a la vista
respondiste una pregunta que cruza hechos, productos (costo) y tiendas
(zona).

\subsection{Índices y planes de ejecución}

Cuando filtras una tabla por una columna (\codigo{WHERE cliente\_id = 5}),
el motor, sin ayuda, tiene que revisar \emph{todas} las filas: un
\termino{escaneo completo} (\emph{full scan}). Un \termino{índice} sobre esa
columna le permite ir directo a las filas que buscas, como el índice
alfabético al final de un libro. \codigo{EXPLAIN QUERY PLAN} te muestra
cómo piensa resolver el motor una consulta, sin ejecutarla.

\begin{rcode}
consulta_cliente <- "SELECT SUM(total) FROM fact_ventas
                     WHERE cliente_id = 5"

# Antes del índice
dbGetQuery(con, paste("EXPLAIN QUERY PLAN", consulta_cliente))$detail

dbExecute(con, "CREATE INDEX idx_fact_cliente
                ON fact_ventas(cliente_id)")

# Después del índice
dbGetQuery(con, paste("EXPLAIN QUERY PLAN", consulta_cliente))$detail
\end{rcode}

\begin{rsalida}
[1] "SCAN fact_ventas"
[1] 0
[1] "SEARCH fact_ventas USING INDEX idx_fact_cliente (cliente_id=?)"
\end{rsalida}

Antes, el plan decía \codigo{SCAN fact\_ventas}: recorrer la tabla completa.
Después, \codigo{SEARCH fact\_ventas USING INDEX idx\_fact\_cliente}: buscar
directamente con el índice. Con 1\,000 filas la diferencia es
imperceptible, pero con 100 millones es la diferencia entre una consulta de
minutos y una de milisegundos.

¿Qué columnas indexar? En un esquema en estrella, las llaves foráneas de la
tabla de hechos (se usan en todos los \emph{joins} y filtros) y las
columnas por las que filtras con frecuencia. Guardamos la creación de
vistas e índices en una función para reutilizarla:

\begin{rcode}
indices <- c(
  "CREATE INDEX IF NOT EXISTS idx_fact_cliente ON fact_ventas(cliente_id)",
  "CREATE INDEX IF NOT EXISTS idx_fact_fecha ON fact_ventas(fecha_id)",
  "CREATE INDEX IF NOT EXISTS idx_fact_producto
     ON fact_ventas(producto_id)",
  "CREATE INDEX IF NOT EXISTS idx_fact_tienda ON fact_ventas(tienda_id)",
  "CREATE INDEX IF NOT EXISTS idx_fact_vendedor
     ON fact_ventas(vendedor_id)"
)

crear_vistas_indices <- function(con) {
  for (sql in c(indices, vistas)) dbExecute(con, sql)
  invisible(TRUE)
}
for (sql in indices) dbExecute(con, sql)

dbGetQuery(con, "SELECT type, name FROM sqlite_master
                 WHERE type IN ('index', 'view') ORDER BY type, name")
\end{rcode}

\begin{rsalida}
   type                         name
1 index             idx_fact_cliente
2 index               idx_fact_fecha
3 index            idx_fact_producto
4 index              idx_fact_tienda
5 index            idx_fact_vendedor
6 index sqlite_autoindex_dim_fecha_1
7  view               vw_kpi_mensual
8  view            vw_ventas_detalle
\end{rsalida}

\codigo{sqlite\_master} es el catálogo interno de SQLite: lista todas las
tablas, vistas e índices. El índice \codigo{sqlite\_autoindex\_\ldots} lo
creó SQLite solo para la restricción \codigo{UNIQUE} de la fecha.
\codigo{IF NOT EXISTS} evita el error si el índice ya existe.

\begin{advertencia}
Los índices no son gratis: ocupan espacio y hacen un poco más lentas las
inserciones, porque cada fila nueva también actualiza cada índice. No
indexes todas las columnas ``por si acaso''; indexa las que realmente usas
para unir y filtrar. En cargas masivas, a veces conviene borrar los
índices, cargar y volver a crearlos.
\end{advertencia}

\begin{hazlotu}
Escribe \codigo{EXPLAIN QUERY PLAN} para una consulta que sume las ventas
de la tienda 3 (\codigo{WHERE tienda\_id = 3}). ¿Usa el índice
\codigo{idx\_fact\_tienda}?
\end{hazlotu}

% ----------------------------------------------------------------------------
\section{Trabajar con datos grandes}
\label{sec:m5-grandes}

Con los datos del curso todo cabe en memoria, pero en tu trabajo te
toparás con tablas de decenas de millones de filas. Tres principios te
ahorrarán muchos problemas:

\begin{enumerate}
  \item \textbf{Agrega en la base, no en R.} Si necesitas ventas por mes,
        pide a la base 12 filas, no 10 millones para agruparlas en R.
  \item \textbf{Trae solo las columnas y filas que necesitas.} Nada de
        \codigo{SELECT *} ni \codigo{dbReadTable()} sobre tablas grandes.
  \item \textbf{Si de verdad necesitas el detalle, tráelo por bloques.}
\end{enumerate}

\begin{ejemplo}{Cuánto pesa traer datos de más}
Compara el tamaño en memoria de tres formas de obtener las ventas
mensuales.
\end{ejemplo}

\begin{rcode}
# a) Todo el detalle, todas las columnas
todo <- dbGetQuery(con, "SELECT * FROM fact_ventas")
# b) Solo las columnas necesarias
dos_columnas <- dbGetQuery(con, "SELECT fecha_id, total FROM fact_ventas")
# c) Ya agregado en la base
agregado <- dbGetQuery(con, "
  SELECT fecha_id / 100 AS anio_mes, SUM(total) AS ventas
  FROM fact_ventas GROUP BY anio_mes")

tibble(forma = c("a) todo", "b) dos columnas", "c) agregado"),
       filas = c(nrow(todo), nrow(dos_columnas), nrow(agregado)),
       kb    = round(c(object.size(todo), object.size(dos_columnas),
                       object.size(agregado)) / 1024, 1))
\end{rcode}

\begin{rsalida}
# A tibble: 3 x 3
  forma           filas    kb
  <chr>           <int> <dbl>
1 a) todo          1000  79.6
2 b) dos columnas  1000  12.6
3 c) agregado        12   1  
\end{rsalida}

Traer toda la tabla ocupa 79.6~KB; solo las dos columnas necesarias, 12.6~KB
(seis veces menos); y el resultado ya agregado, 1~KB con 12 filas. Con
1\,000 filas da igual, pero escala: con 10 millones de filas, la opción a)
ocuparía cientos de megabytes en tu computadora y la c) seguiría siendo de
12 filas. Por cierto, \codigo{fecha\_id / 100} es un truco de la
llave \codigo{AAAAMMDD}: al dividir entre enteros, 20230113 / 100 da 202301,
el año-mes.

Cuando necesitas procesar el detalle completo (por ejemplo, para exportarlo
a otro sistema), usa \codigo{dbSendQuery()} + \codigo{dbFetch(n = \ldots)}.
\codigo{dbSendQuery()} manda la consulta pero no trae datos; cada
\codigo{dbFetch()} trae el siguiente bloque de \codigo{n} filas, y
\codigo{dbHasCompleted()} te dice si ya no quedan más. Así la memoria solo
contiene un bloque a la vez.

\begin{rcode}
resultado <- dbSendQuery(con, "SELECT venta_id, metodo_pago, total
                               FROM fact_ventas")
bloque <- 0
total_acumulado <- 0
while (!dbHasCompleted(resultado)) {
  datos_bloque <- dbFetch(resultado, n = 300)     # 300 filas a la vez
  bloque <- bloque + 1
  total_acumulado <- total_acumulado + sum(datos_bloque$total)
  cat(sprintf("Bloque %d: %d filas, acumulado $%s\n", bloque,
              nrow(datos_bloque),
              format(round(total_acumulado), big.mark = ",")))
}
dbClearResult(resultado)   # liberar el resultado en la base
\end{rcode}

\begin{rsalida}
Bloque 1: 300 filas, acumulado $298,877
Bloque 2: 300 filas, acumulado $584,379
Bloque 3: 300 filas, acumulado $868,613
Bloque 4: 100 filas, acumulado $959,993
\end{rsalida}

Procesamos las 1\,000 ventas en cuatro bloques (tres de 300 y uno de 100),
y el acumulado final coincide con el total del año. Con 10 millones de
filas usarías bloques de 100\,000 y el mismo código funcionaría sin
agotar la memoria. No olvides \codigo{dbClearResult()}: un resultado
pendiente mantiene ocupados recursos en la base.

\begin{consejo}
Para cargas grandes en la otra dirección (de R a la base), escribe por
bloques con \codigo{dbAppendTable()} dentro de una transacción: insertar
un millón de filas en una sola transacción es muchísimo más rápido que en
un millón de transacciones pequeñas. Y si tus datos analíticos crecen
mucho, investiga \paquete{duckdb}: un motor analítico en archivo, tan fácil
de usar como SQLite pero optimizado para agregar millones de filas, que
funciona con \paquete{DBI} y \paquete{dbplyr} exactamente igual.
\end{consejo}

% ----------------------------------------------------------------------------
\section{Conectarte a las bases de datos de tu empresa}
\label{sec:m5-empresa}

Todo lo que aprendiste funciona igual con un servidor real: solo cambia
\codigo{dbConnect()}. El código de esta sección es de referencia (usa
servidores y credenciales ficticias), así que no se ejecuta en el libro.

\subsection{SQL Server (y otros) con ODBC}

\termino{ODBC} (\emph{Open Database Connectivity}) es un estándar para
conectarse a casi cualquier base de datos mediante un \emph{driver}
instalado en el sistema operativo. Es la vía habitual para SQL Server y
Oracle. Primero verifica qué drivers tienes instalados (si no aparece el de
tu motor, pídelo a TI o descárgalo del sitio del fabricante):

\begin{rnoejecutar}
install.packages("odbc")
odbc::odbcListDrivers()          # drivers ODBC instalados en tu equipo
odbc::odbcListDataSources()      # conexiones (DSN) ya configuradas
\end{rnoejecutar}

\begin{rnoejecutar}
# SQL Server con autenticación de Windows (la más común en empresas)
con_sql <- dbConnect(
  odbc::odbc(),
  Driver   = "ODBC Driver 18 for SQL Server",
  Server   = "srv-datos.miempresa.local",
  Database = "VentasDW",
  Trusted_Connection = "yes",
  TrustServerCertificate = "yes"
)

# SQL Server con usuario y contraseña (tomados de variables de entorno)
con_sql <- dbConnect(
  odbc::odbc(),
  Driver   = "ODBC Driver 18 for SQL Server",
  Server   = Sys.getenv("SQLSERVER_HOST"),
  Database = "VentasDW",
  UID      = Sys.getenv("SQLSERVER_USER"),
  PWD      = Sys.getenv("SQLSERVER_PASS")
)

# Desde aquí, todo es igual: dbGetQuery(), tbl(), dbDisconnect()...
# Ojo con el dialecto: SQL Server usa TOP en lugar de LIMIT
dbGetQuery(con_sql, "SELECT TOP 10 * FROM dbo.fact_ventas")
\end{rnoejecutar}

\subsection{PostgreSQL y MySQL/MariaDB}

\begin{rnoejecutar}
install.packages(c("RPostgres", "RMariaDB"))

# PostgreSQL
con_pg <- dbConnect(
  RPostgres::Postgres(),
  host     = Sys.getenv("PG_HOST"),
  port     = 5432,
  dbname   = "analitica",
  user     = Sys.getenv("PG_USER"),
  password = Sys.getenv("PG_PASS")
)

# MySQL o MariaDB
con_my <- dbConnect(
  RMariaDB::MariaDB(),
  host     = Sys.getenv("MYSQL_HOST"),
  port     = 3306,
  dbname   = "tienda_en_linea",
  user     = Sys.getenv("MYSQL_USER"),
  password = Sys.getenv("MYSQL_PASS")
)

dbListTables(con_pg)
dbDisconnect(con_pg)
dbDisconnect(con_my)
\end{rnoejecutar}

\subsection{Credenciales: nunca en el script}

Observa que en ningún ejemplo escribimos la contraseña. Un script con
\codigo{password = "Ventas2024!"} termina tarde o temprano en un correo, en
una carpeta compartida o en un repositorio de Git, y la contraseña queda
expuesta para siempre (borrarla después no basta: Git guarda el historial).
La práctica estándar es guardar las credenciales en
\termino{variables de entorno}, en el archivo \archivo{.Renviron} de tu
carpeta de usuario, y leerlas con \codigo{Sys.getenv()}.

Abre el archivo con \codigo{usethis::edit\_r\_environ()} (o créalo a mano)
y escribe una línea por variable, sin comillas ni espacios alrededor del
\codigo{=}; los valores de este ejemplo son ficticios:

\begin{textocodigo}[title=Archivo .Renviron (ejemplo ficticio)]
SQLSERVER_HOST=srv-datos.miempresa.local
SQLSERVER_USER=analista_bi
SQLSERVER_PASS=cambia-esta-clave
PG_HOST=pg.miempresa.local
PG_USER=lector_bi
PG_PASS=otra-clave-ficticia
\end{textocodigo}

Reinicia R (\menu{Session > Restart R}) para que lea el archivo. A partir de
ahí, \codigo{Sys.getenv("SQLSERVER\_USER")} devuelve
\codigo{"analista\_bi"} y tu script no contiene ningún secreto:

\begin{rnoejecutar}
usethis::edit_r_environ()          # abre ~/.Renviron para editarlo
Sys.getenv("SQLSERVER_USER")       # "analista_bi"
nchar(Sys.getenv("SQLSERVER_PASS")) > 0   # comprobar sin imprimirla
\end{rnoejecutar}

\begin{advertencia}
\begin{itemize}
  \item Nunca subas \archivo{.Renviron} a Git: agrégalo a tu
        \archivo{.gitignore}.
  \item No imprimas la contraseña en la consola ni en un reporte.
  \item En RStudio también puedes pedir la contraseña al momento con
        \codigo{rstudioapi::askForPassword("Contraseña")}, útil en scripts
        interactivos.
\end{itemize}
\end{advertencia}

\begin{negocio}{Permisos y buenas prácticas en la empresa}
\begin{itemize}
  \item Pide a TI un usuario de \textbf{solo lectura} para analizar. Así es
        imposible que un \codigo{DELETE} accidental borre datos de
        producción.
  \item Consulta de preferencia una \textbf{réplica} o el \emph{data
        warehouse}, no la base transaccional que usa el punto de venta.
  \item Evita consultas pesadas en horario pico, y usa \codigo{LIMIT}/
        \codigo{TOP} mientras exploras.
  \item En aplicaciones Shiny con muchos usuarios no abras una conexión por
        cada consulta: usa el paquete \paquete{pool}, que mantiene un
        conjunto de conexiones abiertas y las reparte
        (\codigo{pool::dbPool()} en lugar de \codigo{dbConnect()}). Lo
        retomarás en el Módulo~\ref{cap:shiny}.
\end{itemize}
\end{negocio}

% ----------------------------------------------------------------------------
\section{Caso práctico integral: el data warehouse de ventas}
\label{sec:m5-caso}

\begin{ejemplo}{Del CSV al tablero, pasando por el data warehouse}
La dirección te pide dejar lista una base analítica que cualquiera pueda
consultar, y responder con ella siete preguntas para la junta mensual:
\begin{enumerate}
  \item ¿Cómo evolucionaron las ventas y el margen mes a mes?
  \item ¿Cuáles son los productos que más venden y cuánto margen dejan?
  \item ¿Quiénes son los clientes más valiosos?
  \item ¿Qué tiendas son más productivas por metro cuadrado?
  \item ¿Quiénes son los dos mejores vendedores de cada zona?
  \item ¿Cuánto crecemos mes contra mes?
  \item ¿Se compra distinto en fin de semana?
\end{enumerate}
\end{ejemplo}

\subsection{Paso 1: construir la base desde cero}

Gracias a que escribimos todo como funciones, reconstruir el \emph{data
warehouse} completo es un script corto. Cerramos la conexión actual,
borramos el archivo y ejecutamos cada etapa. Esta vez la base solo
contendrá el esquema en estrella (sin las tablas crudas del principio): es
lo que publicarías para los usuarios de negocio.

\begin{rcode}
dbDisconnect(con)

reconstruir_dw <- function(ruta_bd, archivos_ventas) {
  unlink(ruta_bd)                                 # empezar desde cero
  con <- dbConnect(RSQLite::SQLite(), ruta_bd)
  crear_esquema(con)                              # tablas y llaves
  cargar_dimensiones(con, preparar_dimensiones()) # dimensiones
  for (archivo in archivos_ventas) {              # hechos (ETL)
    ejecutar_etl(con, archivo)
  }
  crear_vistas_indices(con)                       # vistas e índices
  con
}

con <- reconstruir_dw("datos/bi_retail.sqlite",
                      c("datos/entrada/ventas_2023_s1.csv",
                        "datos/entrada/ventas_2023_s2.csv"))
dbListTables(con)
\end{rcode}

\begin{rsalida}
ventas_2023_s1.csv -> OK: 511 leídas, 511 nuevas 
ventas_2023_s2.csv -> OK: 661 leídas, 489 nuevas 
 [1] "dim_cliente"       "dim_fecha"         "dim_producto"     
 [4] "dim_tienda"        "dim_vendedor"      "fact_ventas"      
 [7] "log_etl"           "sqlite_sequence"   "vw_kpi_mensual"   
[10] "vw_ventas_detalle"
\end{rsalida}

En segundos tienes la base completa: cinco dimensiones, la tabla de hechos
con sus 1\,000 ventas, la bitácora y dos vistas. Si mañana llega un archivo
nuevo, solo hay que llamar a \codigo{ejecutar\_etl()}.

\subsection{Paso 2: responder las preguntas con SQL}

\subsubsection{Pregunta 1. Ventas y margen por mes}

\begin{rcode}
kpi_mensual <- dbGetQuery(con, "
  SELECT mes, nombre_mes, tickets, ventas, margen,
         ROUND(100.0 * margen / ventas, 1) AS margen_pct
  FROM vw_kpi_mensual
  ORDER BY anio, mes
")
kpi_mensual
\end{rcode}

\begin{rsalida}
   mes nombre_mes tickets   ventas   margen margen_pct
1    1      enero      98 82559.66 34415.12       41.7
2    2    febrero      79 65899.34 29212.36       44.3
3    3      marzo      91 90930.61 46572.20       51.2
4    4      abril      71 57037.74 19817.86       34.7
5    5       mayo      88 93910.21 49550.60       52.8
6    6      junio      84 88509.12 42049.19       47.5
7    7      julio      83 86456.05 41120.26       47.6
8    8     agosto      75 78929.49 33268.92       42.2
9    9 septiembre      86 68855.87 29467.75       42.8
10  10    octubre      74 70004.10 31690.67       45.3
11  11  noviembre      80 80506.89 36578.96       45.4
12  12  diciembre      91 96393.61 50893.04       52.8
\end{rsalida}

Las ventas oscilan entre \$57\,038 (abril) y \$96\,394 (diciembre), y el
margen porcentual va de 34.7\% en abril a 52.8\% en mayo y diciembre. Abril
fue malo por partida doble: vendió menos \emph{y} con menor margen. ¿Podemos confiar en estas cifras? Una buena práctica es
\emph{reconciliar}: calcular lo mismo por otro camino y comprobar que
coincide. Hagámoslo con \paquete{dplyr} directamente sobre el CSV original,
sin pasar por la base:

\begin{rcode}
kpi_dplyr <- read_csv("datasets/transacciones.csv",
                      show_col_types = FALSE) %>%
  group_by(mes = month(fecha)) %>%
  summarise(tickets = n(),
            ventas  = round(sum(cantidad * precio_venta), 2),
            .groups = "drop")

comparacion <- kpi_mensual %>%
  select(mes, tickets, ventas) %>%
  inner_join(kpi_dplyr, by = "mes", suffix = c("_sql", "_dplyr"))

all.equal(comparacion$ventas_sql, comparacion$ventas_dplyr)
all(comparacion$tickets_sql == comparacion$tickets_dplyr)
\end{rcode}

\begin{rsalida}
[1] TRUE
[1] TRUE
\end{rsalida}

Las ventas y los tickets de los 12 meses coinciden exactamente: el ETL no
perdió ni duplicó nada. Con los datos ya validados, la gráfica para la
junta es inmediata (figura~\ref{fig:m05-ventas-mensuales}):

\begin{rcode}
color_principal <- "#2a78d6"

grafica_mensual <- kpi_mensual %>%
  mutate(nombre_mes = factor(nombre_mes, levels = nombre_mes)) %>%
  ggplot(aes(x = nombre_mes, y = ventas, group = 1)) +
  geom_line(color = color_principal, linewidth = 0.8) +
  geom_point(color = color_principal, size = 2) +
  scale_y_continuous(labels = scales::dollar,
                     limits = c(0, NA)) +
  labs(title = "Ventas mensuales 2023",
       subtitle = "Fuente: vista vw_kpi_mensual del data warehouse",
       x = NULL, y = "Ventas") +
  theme_minimal() +
  theme(panel.grid.minor = element_blank(),
        axis.text.x = element_text(angle = 45, hjust = 1))
grafica_mensual
\end{rcode}

\begin{figure}[H]
\centering
\includegraphics[width=0.9\textwidth]{m05-ventas-mensuales}
\caption{Ventas mensuales calculadas con SQL a partir del data warehouse.}
\label{fig:m05-ventas-mensuales}
\end{figure}

\subsubsection{Pregunta 2. Productos que más venden y su margen}

\begin{rcode}
dbGetQuery(con, "
  SELECT p.nombre_producto, p.categoria,
         SUM(f.cantidad)                            AS unidades,
         ROUND(SUM(f.total), 0)                     AS ventas,
         ROUND(SUM(f.margen), 0)                    AS margen,
         ROUND(100.0 * SUM(f.margen) / SUM(f.total), 1) AS margen_pct
  FROM fact_ventas AS f
  JOIN dim_producto AS p ON f.producto_id = p.producto_id
  GROUP BY p.producto_id
  ORDER BY ventas DESC
  LIMIT 5
")
\end{rcode}

\begin{rsalida}
  nombre_producto  categoria unidades ventas margen margen_pct
1     Producto 06   Juguetes       82  34028  27095       79.6
2     Producto 19   Deportes       85  28368   7557       26.6
3     Producto 31  Alimentos       74  25430  23622       92.9
4     Producto 47 Automotriz       82  24527  20017       81.6
5     Producto 46 Automotriz       74  24470  11333       46.3
\end{rsalida}

El Producto 06 (Juguetes) lidera en ventas y además deja un margen de
79.6\%. En cambio, el Producto 19 (Deportes), segundo en ventas, apenas
deja 26.6\%: vende mucho pero aporta poco. Un producto que vende mucho no
necesariamente es el más rentable.

\begin{rcode}
# ¿Hay productos que se venden con pérdida?
dbGetQuery(con, "
  SELECT p.nombre_producto, p.costo,
         ROUND(AVG(f.precio_venta), 2) AS precio_venta_prom,
         ROUND(SUM(f.margen), 0)       AS margen_total
  FROM fact_ventas AS f
  JOIN dim_producto AS p ON f.producto_id = p.producto_id
  GROUP BY p.producto_id
  HAVING SUM(f.margen) < 0
  ORDER BY margen_total
")
\end{rcode}

\begin{rsalida}
  nombre_producto  costo precio_venta_prom margen_total
1     Producto 13 392.63            257.58        -8602
2     Producto 35 353.50            230.62        -6530
3     Producto 41 389.86            294.97        -6046
4     Producto 12 368.62            296.11        -4332
5     Producto 02 399.12            304.00        -4137
6     Producto 29 391.67            355.31        -2258
7     Producto 10 283.26            260.12        -1282
\end{rsalida}

Siete productos se vendieron con pérdida en el año: su costo unitario es
mayor que el precio promedio al que se vendieron. El Producto 13, por
ejemplo, cuesta \$392.63 y se vendió en promedio a \$257.58, con una
pérdida acumulada de \$8\,602. Es el hallazgo de calidad de datos que ya
conocías (precios de venta generados sin relación con el costo), pero en
una empresa real sería una alerta roja para el área comercial.

\subsubsection{Pregunta 3. Clientes más valiosos}

\begin{rcode}
dbGetQuery(con, "
  SELECT c.nombre, c.ciudad, c.segmento,
         CASE c.es_premium WHEN 1 THEN 'Sí' ELSE 'No' END AS premium,
         COUNT(*)               AS compras,
         ROUND(SUM(f.total), 0) AS gasto,
         ROUND(SUM(f.margen), 0) AS margen
  FROM fact_ventas AS f
  JOIN dim_cliente AS c ON f.cliente_id = c.cliente_id
  GROUP BY c.cliente_id
  ORDER BY gasto DESC
  LIMIT 5
")
\end{rcode}

\begin{rsalida}
          nombre      ciudad segmento premium compras gasto margen
1     Pedro Cruz      Toluca   Adulto      No      13 14680   5882
2  Pedro Morales        León   Senior      No       9 12504   8669
3 Carmen Jiménez Guadalajara   Adulto      No       9 11848   5001
4 Sofia González      Cancún   Maduro      No      12 11676   6778
5  Alberto Reyes Guadalajara   Maduro      No      13 11356   5905
\end{rsalida}

Pedro Cruz, de Toluca, es el cliente más valioso por gasto (\$14\,680 en 13
compras); es el cliente 5 que ya habíamos visto. Pero si ordenamos por
margen, Pedro Morales, de León, deja más (\$8\,669 con solo 9 compras). Un
dato curioso: ninguno de los cinco es cliente premium, una pista para
revisar los criterios del programa.

\subsubsection{Pregunta 4. Productividad de las tiendas}

\begin{rcode}
dbGetQuery(con, "
  SELECT t.nombre_tienda, t.tipo, t.tamano_m2,
         ROUND(SUM(f.total), 0)               AS ventas,
         ROUND(SUM(f.total) / t.tamano_m2, 1) AS ventas_m2,
         ROUND(AVG(f.total), 0)               AS ticket_prom
  FROM fact_ventas AS f
  JOIN dim_tienda AS t ON f.tienda_id = t.tienda_id
  GROUP BY t.tienda_id
  ORDER BY ventas_m2 DESC
")
\end{rcode}

\begin{rsalida}
      nombre_tienda      tipo tamano_m2 ventas ventas_m2 ticket_prom
1  Sucursal Express   Express       300 102791     342.6         952
2   Sucursal Outlet    Outlet       350 114377     326.8        1079
3     Sucursal Este  Estándar       400  88420     221.0         941
4    Sucursal Norte  Estándar       450  89024     197.8         899
5   Sucursal Centro Principal       500  88414     176.8        1016
6      Sucursal Sur  Estándar       600 105305     175.5         966
7  Sucursal Plaza B  Estándar       480  80236     167.2         944
8    Sucursal Oeste  Estándar       550  84672     153.9         891
9  Sucursal Plaza A Principal       700 105216     150.3         965
10 Sucursal Premium   Premium       800 101538     126.9         940
\end{rsalida}

Por metro cuadrado, las tiendas pequeñas ganan: la Sucursal Express
(300~m²) vende \$342.6 por m² y el Outlet \$326.8, mientras que la Sucursal
Premium, la más grande (800~m²), queda al final con \$126.9. Las ventas
totales son parecidas entre tiendas, así que el tamaño no se traduce en más
ventas: una conclusión valiosa para decidir el tamaño de las próximas
aperturas.

\subsubsection{Pregunta 5. Los dos mejores vendedores de cada zona}

Aquí combinamos una CTE con una función de ventana particionada por zona.
Un vendedor puede vender en tiendas de varias zonas, así que el ranking es
``por lo que vendió en esa zona''.

\begin{rcode}
dbGetQuery(con, "
  WITH ventas_vendedor_zona AS (
    SELECT t.zona, v.nombre AS vendedor, SUM(f.total) AS ventas
    FROM fact_ventas AS f
    JOIN dim_tienda   AS t ON f.tienda_id   = t.tienda_id
    JOIN dim_vendedor AS v ON f.vendedor_id = v.vendedor_id
    GROUP BY t.zona, v.vendedor_id
  ),
  ranking AS (
    SELECT zona, vendedor, ROUND(ventas, 0) AS ventas,
           RANK() OVER (PARTITION BY zona ORDER BY ventas DESC) AS lugar
    FROM ventas_vendedor_zona
  )
  SELECT * FROM ranking
  WHERE lugar <= 2
  ORDER BY zona, lugar
")
\end{rcode}

\begin{rsalida}
     zona        vendedor ventas lugar
1  Centro   Isabel García  21198     1
2  Centro    Sofia Rivera  17066     2
3    Este  Fernando Pérez   8881     1
4    Este   Pedro Morales   7589     2
5   Norte Ricardo Ramírez  22238     1
6   Norte  Marta González  21909     2
7   Oeste   Isabel García  11709     1
8   Oeste    Juan Sánchez  10093     2
9     Sur     Marta Gómez  11526     1
10    Sur    Marta Rivera   9124     2
\end{rsalida}

Ricardo Ramírez lidera la zona Norte (\$22\,238) y Isabel García aparece
como la mejor vendedora \emph{en dos zonas} (Centro y Oeste), señal de que
atiende en varias sucursales. En la zona Este, con una sola tienda, las
cifras son mucho menores. Recuerda el hallazgo anterior: varios de estos
``vendedores'' no pertenecen al departamento de Ventas.

\subsubsection{Pregunta 6. Crecimiento mes contra mes}

\begin{rcode}
dbGetQuery(con, "
  SELECT nombre_mes, ventas,
         ROUND(100.0 * (ventas - LAG(ventas) OVER (ORDER BY mes)) /
               LAG(ventas) OVER (ORDER BY mes), 1) AS crec_ventas_pct,
         ROUND(100.0 * (margen - LAG(margen) OVER (ORDER BY mes)) /
               LAG(margen) OVER (ORDER BY mes), 1) AS crec_margen_pct
  FROM vw_kpi_mensual
  ORDER BY mes
")
\end{rcode}

\begin{rsalida}
   nombre_mes   ventas crec_ventas_pct crec_margen_pct
1       enero 82559.66              NA              NA
2     febrero 65899.34           -20.2           -15.1
3       marzo 90930.61            38.0            59.4
4       abril 57037.74           -37.3           -57.4
5        mayo 93910.21            64.6           150.0
6       junio 88509.12            -5.8           -15.1
7       julio 86456.05            -2.3            -2.2
8      agosto 78929.49            -8.7           -19.1
9  septiembre 68855.87           -12.8           -11.4
10    octubre 70004.10             1.7             7.5
11  noviembre 80506.89            15.0            15.4
12  diciembre 96393.61            19.7            39.1
\end{rsalida}

El margen se mueve mucho más que las ventas: en mayo las ventas crecieron
64.6\% pero el margen 150\%, y en abril las ventas cayeron 37.3\% y el
margen 57.4\%. Es decir, la mezcla de productos vendidos cambia de un mes a
otro, y los meses buenos lo son también por vender productos más
rentables.

\subsubsection{Pregunta 7. Fin de semana contra entre semana}

\begin{rcode}
dbGetQuery(con, "
  SELECT CASE d.es_fin_semana WHEN 1 THEN 'Fin de semana'
                              ELSE 'Entre semana' END AS tipo_dia,
         COUNT(DISTINCT d.fecha_id)                    AS dias,
         COUNT(*)                                      AS tickets,
         ROUND(1.0 * COUNT(*) / COUNT(DISTINCT d.fecha_id), 2)
                                                  AS tickets_por_dia,
         ROUND(AVG(f.total), 2)                        AS ticket_prom
  FROM fact_ventas AS f
  JOIN dim_fecha AS d ON f.fecha_id = d.fecha_id
  GROUP BY tipo_dia
")
\end{rcode}

\begin{rsalida}
       tipo_dia dias tickets tickets_por_dia ticket_prom
1  Entre semana  246     724            2.94      961.49
2 Fin de semana   99     276            2.79      956.07
\end{rsalida}

Hubo ventas en 246 días entre semana y 99 de fin de semana (345 días en
total: 20 días del año no registraron ninguna venta). Entre semana se
venden 2.94 tickets por día contra 2.79 en fin de semana, y el ticket
promedio es prácticamente igual (\$961 contra \$956). No hay un patrón de
fin de semana marcado. Este tipo de pregunta es trivial gracias a la dimensión de
fechas: el atributo \codigo{es\_fin\_semana} ya está calculado y cualquier
analista puede usarlo sin saber manipular fechas.

\begin{negocio}{Qué entregarías en la junta}
Con estas siete consultas tienes el material de un reporte ejecutivo
completo, y la base \archivo{bi\_retail.sqlite} queda como la fuente única
de la verdad: el Módulo~\ref{cap:shiny} construirá un tablero que lee de
ella y el Módulo~\ref{cap:reportes} un reporte automático. Lo más valioso
no son las cifras de hoy, sino que el próximo mes basta ejecutar
\codigo{ejecutar\_etl()} con el archivo nuevo y volver a correr las
consultas.
\end{negocio}

\begin{hazlotu}
Agrega una octava pregunta: ¿qué método de pago predomina en cada zona?
Usa \codigo{vw\_ventas\_detalle} y \codigo{ROW\_NUMBER()} particionado por
zona.
\end{hazlotu}

% ----------------------------------------------------------------------------
\section{Errores comunes}
\label{sec:m5-errores}

Para ver los mensajes reales sin detener el script, usamos esta pequeña
función que ejecuta una expresión y, si falla, imprime el mensaje de error
(o de advertencia) en lugar de detenerse:

\begin{rcode}
mostrar_error <- function(expr) {
  tryCatch(expr,
           error   = function(e) cat("Error:", conditionMessage(e), "\n"),
           warning = function(w) cat("Aviso:", conditionMessage(w), "\n"))
}
\end{rcode}

\subsubsection{``no such table'' y ``no such column''}

\begin{rcode}
mostrar_error(dbGetQuery(con, "SELECT * FROM ventas_2024"))
mostrar_error(dbGetQuery(con, "SELECT totl FROM fact_ventas"))
\end{rcode}

\begin{rsalida}
Error: no such table: ventas_2024 
Error: no such column: totl 
\end{rsalida}

La tabla o la columna no existe, casi siempre por un error de dedo o
porque estás conectado a otra base. Revisa con \codigo{dbListTables(con)} y
\codigo{dbListFields(con, "tabla")}. Recuerda también que, tras reconstruir
la base, las tablas crudas (\codigo{transacciones}, \codigo{clientes}\ldots)
ya no existen: solo el esquema en estrella.

\subsubsection{La tabla ya existe}

\begin{rcode}
mostrar_error(dbWriteTable(con, "dim_tienda", dimensiones$dim_tienda))
\end{rcode}

\begin{rsalida}
Error: Table dim_tienda exists in database, and both overwrite and append are FALSE 
\end{rsalida}

\codigo{dbWriteTable()} se niega a pisar una tabla existente. Decide qué
quieres: \codigo{overwrite = TRUE} la reemplaza (¡perdiendo llaves y
reglas si la habías creado con \codigo{CREATE TABLE}!) y
\codigo{append = TRUE} agrega filas. Para tablas diseñadas por ti, usa
\codigo{dbAppendTable()}.

\subsubsection{Llaves duplicadas al volver a cargar}

\begin{rcode}
mostrar_error(dbAppendTable(con, "dim_tienda", dimensiones$dim_tienda))
\end{rcode}

\begin{rsalida}
Error: UNIQUE constraint failed: dim_tienda.tienda_id 
\end{rsalida}

Intentaste cargar filas cuya llave primaria ya existe (típico al correr dos
veces un script). Usa carga incremental (\codigo{anti\_join()} o
\codigo{INSERT OR IGNORE}) o borra primero lo que vas a recargar.

\subsubsection{Fechas convertidas en números}

Si ves fechas como \codigo{19370} o años como \codigo{-4672}, escribiste un
\codigo{Date} de R directamente en SQLite (sección~\ref{sec:m5-fechas}).
Convierte las fechas a texto con \codigo{format(fecha, "\%Y-\%m-\%d")}
antes de escribirlas y con \codigo{as.Date()} al leerlas.

\subsubsection{Comillas simples contra comillas dobles}

En SQL estándar, el texto va entre comillas \emph{simples}
(\codigo{'CDMX'}) y las comillas \emph{dobles} son para nombres de
columnas o tablas. SQLite, por compatibilidad, acepta comillas dobles para
texto si no encuentra una columna con ese nombre, lo que produce errores
silenciosos:

\begin{rcode}
dbGetQuery(con, "SELECT COUNT(*) AS n FROM dim_cliente
                 WHERE ciudad = 'CDMX'")     # correcto: 27
dbGetQuery(con, "SELECT COUNT(*) AS n FROM dim_cliente
                 WHERE ciudad = \"ciudad\"")  # ¡columna consigo misma!
\end{rcode}

\begin{rsalida}
   n
1 27
    n
1 200
\end{rsalida}

La segunda consulta pretendía buscar clientes de una ciudad llamada
``ciudad'', pero \codigo{"ciudad"} se interpretó como la \emph{columna}
ciudad, que siempre es igual a sí misma: devolvió los 200 clientes. En
PostgreSQL o SQL Server, \codigo{"CDMX"} directamente marcaría error de
columna inexistente. Usa siempre comillas simples para texto (y
\codigo{params} para valores variables).

\subsubsection{Olvidar cerrar conexiones y resultados}

\begin{rcode}
con_temporal <- dbConnect(RSQLite::SQLite(), ":memory:")
dbDisconnect(con_temporal)
mostrar_error(dbGetQuery(con_temporal, "SELECT 1"))
\end{rcode}

\begin{rsalida}
Error: Invalid or closed connection 
\end{rsalida}

Usar una conexión ya cerrada da \codigo{Invalid or closed connection}: vuelve
a conectarte. El error contrario (nunca cerrar) no da error inmediato, pero
R te avisa más tarde con un mensaje como \emph{call dbDisconnect() when
finished working with a connection}, y en un servidor acumula conexiones
abiertas. Usa \codigo{dbDisconnect()} al final de cada script y
\codigo{on.exit()} dentro de las funciones. Del mismo modo, cada
\codigo{dbSendQuery()} necesita su \codigo{dbClearResult()}.

\begin{table}[H]
\centering
\small
\caption{Resumen de errores frecuentes con bases de datos.}
\label{tab:m5-errores}
\begin{tabularx}{\textwidth}{@{}>{\raggedright\arraybackslash}p{5.2cm}X@{}}
\toprule
\textbf{Mensaje o síntoma} & \textbf{Causa y solución} \\
\midrule
\codigo{no such table: X} & La tabla no existe o estás en otra base.
Revisa \codigo{dbListTables()}. \\
\codigo{no such column: X} & Error de dedo o columna de otra tabla; revisa
\codigo{dbListFields()} y los alias. \\
\codigo{ambiguous column name: X} & La columna existe en dos tablas del
JOIN: antepón el alias (\codigo{t.producto\_id}). \\
\codigo{Table X exists in database\ldots} & Usa \codigo{overwrite},
\codigo{append} o \codigo{dbAppendTable()}. \\
\codigo{UNIQUE constraint failed} & Llave repetida: carga incremental o
\codigo{INSERT OR IGNORE}. \\
\codigo{FOREIGN KEY constraint failed} & Un valor no existe en la
dimensión: carga primero las dimensiones o corrige el dato. \\
Fechas como números & Convierte \codigo{Date} a texto ISO antes de
escribir. \\
\codigo{Invalid or closed connection} & La conexión se cerró: vuelve a
conectarte. \\
Resultados ``demasiado buenos'' & Comillas dobles en lugar de simples, o
un JOIN que multiplica filas. \\
\bottomrule
\end{tabularx}
\end{table}

Terminamos cerrando la conexión principal del módulo:

\begin{rcode}
dbDisconnect(con)
\end{rcode}

% ----------------------------------------------------------------------------
\begin{resumen}
\begin{itemize}
  \item Las empresas guardan sus datos en bases de datos por volumen,
        integridad, concurrencia y seguridad. Las OLTP operan el negocio;
        los \emph{data warehouses} (OLAP) sirven para analizarlo.
  \item \paquete{DBI} ofrece las mismas funciones para cualquier motor:
        \codigo{dbConnect()}, \codigo{dbWriteTable()},
        \codigo{dbGetQuery()}, \codigo{dbExecute()},
        \codigo{dbDisconnect()}. Solo cambia el driver.
  \item En SQLite guarda las fechas como texto \codigo{'AAAA-MM-DD'}; si no,
        se vuelven números.
  \item SQL: \codigo{SELECT}, \codigo{WHERE}, \codigo{GROUP BY},
        \codigo{HAVING}, \codigo{ORDER BY}, \codigo{JOIN}, subconsultas,
        CTE y funciones de ventana responden casi cualquier pregunta.
  \item Nunca pegues valores del usuario con \codigo{paste()}: usa
        \codigo{params}, \codigo{glue\_sql()} o \codigo{sqlInterpolate()}.
  \item \paquete{dbplyr} traduce dplyr a SQL; \codigo{show\_query()} te
        enseña el SQL y \codigo{collect()} trae el resultado.
  \item Un data warehouse en estrella tiene una tabla de hechos y
        dimensiones con llaves declaradas, vistas para los KPIs e índices
        en las llaves.
  \item Un ETL confiable es idempotente (sin duplicados), atómico
        (transacciones) y trazable (bitácora).
  \item Con datos grandes: agrega en la base, trae solo lo necesario y usa
        \codigo{dbFetch()} por bloques.
  \item Credenciales en \archivo{.Renviron}, nunca en el script; usuarios de
        solo lectura para analizar.
\end{itemize}

\medskip
\centering
\small
\begin{tabular}{@{}ll@{}}
\toprule
\textbf{SQL} & \textbf{dplyr} \\
\midrule
\codigo{SELECT col1, col2} & \codigo{select(col1, col2)} \\
\codigo{SELECT expr AS nueva} & \codigo{mutate(nueva = expr)} \\
\codigo{WHERE cond} & \codigo{filter(cond)} \\
\codigo{ORDER BY col DESC} & \codigo{arrange(desc(col))} \\
\codigo{LIMIT 5} & \codigo{head(5)} / \codigo{slice\_head(n = 5)} \\
\codigo{GROUP BY g} + \codigo{SUM(x)} & \codigo{group\_by(g)} +
  \codigo{summarise(sum(x))} \\
\codigo{HAVING cond} & \codigo{filter()} después de \codigo{summarise()} \\
\codigo{COUNT(DISTINCT x)} & \codigo{n\_distinct(x)} \\
\codigo{CASE WHEN \ldots{} END} & \codigo{case\_when()} \\
\codigo{INNER JOIN} / \codigo{LEFT JOIN} & \codigo{inner\_join()} /
  \codigo{left\_join()} \\
\codigo{LEFT JOIN \ldots{} WHERE b.id IS NULL} & \codigo{anti\_join()} \\
\codigo{WITH paso AS (\ldots)} & guardar un resultado intermedio \\
\codigo{ROW\_NUMBER() OVER (PARTITION BY g \ldots)} & \codigo{group\_by(g)}
  + \codigo{row\_number()} \\
\codigo{SUM(x) OVER (ORDER BY t)} & \codigo{cumsum(x)} \\
\codigo{LAG(x) OVER (ORDER BY t)} & \codigo{lag(x)} \\
\bottomrule
\end{tabular}
\end{resumen}

% ----------------------------------------------------------------------------
\section{Ejercicios}
\label{sec:m5-ejercicios}

Resuelve cada ejercicio en un script nuevo. Para no depender de la base del
capítulo, trabaja con una base en memoria
(\codigo{dbConnect(RSQLite::SQLite(), ":memory:")}) y carga solo las tablas
que necesites. Salvo que se indique otra cosa, responde con SQL escrito a
mano mediante \codigo{dbGetQuery()}.

\begin{ejercicio}{m5-1}{Tu primera base de datos}
\nivel{Básico}
Crea una base en memoria, carga \archivo{productos.csv} en una tabla
llamada \codigo{productos} y muestra sus tablas y columnas con
\codigo{dbListTables()} y \codigo{dbListFields()}. Después, con SQL,
calcula cuántos productos hay por categoría y el precio de catálogo
promedio, ordenado de más a menos productos. Cierra la conexión al final.
\end{ejercicio}

\begin{ejercicio}{m5-2}{Alerta de inventario}
\nivel{Básico}
El área de compras necesita la lista de productos \emph{activos} cuyo
\codigo{stock\_actual} está por debajo del \codigo{stock\_minimo}. Muestra
nombre, categoría, proveedor, stock actual, stock mínimo y las piezas
faltantes (\codigo{stock\_minimo - stock\_actual}), ordenado de mayor a
menor faltante. ¿Cuántos productos están en alerta y qué proveedor
aparece más?
\end{ejercicio}

\begin{ejercicio}{m5-3}{Nómina por departamento}
\nivel{Básico}
Con \archivo{empleados.csv}, calcula por departamento: número de
empleados, salario mensual promedio y salario máximo. Muestra solo los
departamentos con al menos 15 empleados (\codigo{HAVING}), ordenados por
salario promedio descendente.
\end{ejercicio}

\begin{ejercicio}{m5-4}{Ventas por ciudad y productos sin ventas en diciembre}
\nivel{Intermedio}
Carga \archivo{transacciones.csv}, \archivo{tiendas.csv} y
\archivo{productos.csv} (con fechas como texto ISO). (a) Con un
\codigo{INNER JOIN}, calcula por ciudad de la tienda el número de tickets,
las ventas y el ticket promedio. (b) Con un \codigo{LEFT JOIN}, encuentra
los productos que no tuvieron ninguna venta en diciembre de 2023.
\end{ejercicio}

\begin{ejercicio}{m5-5}{Fechas bien guardadas}
\nivel{Intermedio}
Carga \archivo{ventas\_retail.csv} dos veces en una base en memoria: una
tabla \codigo{ventas\_mal} con la fecha tal como la lee
\codigo{read\_csv()} y otra \codigo{ventas} con la fecha como texto ISO.
Demuestra con \codigo{typeof()} y \codigo{strftime()} la diferencia. Luego,
con la tabla correcta, calcula por trimestre (a partir de la fecha) las
ventas totales, el descuento promedio y el número de días con descuento.
\end{ejercicio}

\begin{ejercicio}{m5-6}{Consulta segura de un cliente}
\nivel{Intermedio}
Escribe una función \codigo{ventas\_cliente(con, cliente, desde, hasta)}
que devuelva el número de compras y el gasto de un cliente en un periodo,
usando una consulta parametrizada con marcadores con nombre. Pruébala con
el cliente 5 en el primer semestre de 2023 y en el segundo. Después
demuestra que el texto \codigo{"5 OR 1=1"} no logra traer las compras de
todos los clientes.
\end{ejercicio}

\begin{ejercicio}{m5-7}{Lo mismo con dbplyr}
\nivel{Intermedio}
Con \paquete{dbplyr}, calcula para cada método de pago y trimestre el
número de tickets y las ventas, a partir de la tabla \codigo{transacciones}.
Muestra el SQL generado con \codigo{show\_query()}, trae el resultado con
\codigo{collect()} y comprueba con \codigo{all.equal()} que coincide con
la misma consulta escrita a mano en SQL.
\end{ejercicio}

\begin{ejercicio}{m5-8}{Rankings y Pareto con funciones de ventana}
\nivel{Avanzado}
(a) Encuentra los dos clientes con mayor gasto de cada ciudad (ciudad del
cliente) usando \codigo{RANK()} con \codigo{PARTITION BY}; muestra solo
Monterrey, Guadalajara y CDMX. (b) Construye la curva de Pareto de
productos: ordena los productos por ventas y calcula con
\codigo{SUM() OVER} el porcentaje acumulado. ¿Cuántos productos generan el
80\% de las ventas?
\end{ejercicio}

\begin{ejercicio}{m5-9}{Una estrella en miniatura}
\nivel{Avanzado}
A partir de \archivo{ventas\_retail.csv}, \archivo{productos.csv} y
\archivo{tiendas.csv}, crea en una base en memoria un esquema en estrella
pequeño: \codigo{dim\_producto}, \codigo{dim\_tienda} y
\codigo{fact\_ventas\_diarias} (llave primaria \codigo{fecha}, llaves
foráneas a ambas dimensiones, \codigo{CHECK (cantidad > 0)}). Carga los
datos con \codigo{dbAppendTable()}, crea una vista con las ventas por
trimestre y tipo de tienda, y muestra con \codigo{EXPLAIN QUERY PLAN} el
efecto de un índice sobre \codigo{tienda\_id}.
\end{ejercicio}

\begin{ejercicio}{m5-10}{ETL incremental con bitácora}
\nivel{Reto}
Construye un ETL para \archivo{ventas\_retail.csv} en una base en memoria:
una tabla \codigo{fact\_ventas\_diarias} con llave primaria \codigo{fecha}
(texto) y una tabla \codigo{log\_etl}. Escribe las funciones
\codigo{extraer\_ventas()}, \codigo{transformar\_ventas()},
\codigo{cargar\_ventas()} (que inserte desde una tabla temporal solo las
fechas que aún no existen) y \codigo{ejecutar\_etl()} (con transacción y registro en la
bitácora, también de los errores). Simula tres lotes: enero-junio,
abril-septiembre (traslapado) y octubre-diciembre; ejecuta además un lote
repetido y un lote con una cantidad negativa que debe ser rechazado por la
regla \codigo{CHECK}. Muestra la bitácora y verifica que la tabla termina
con exactamente 365 filas y el mismo total de ventas que el CSV.
\end{ejercicio}
